% =============================================================================
%  progress2.tex — รายงานความก้าวหน้าโครงการ ULMs ครั้งที่ 2
%
%  build:  cd ~/Projects/SoftwareEn/docs && latexmk progress2.tex
%          ผลลัพธ์ออกที่ build/progress2.pdf (ตาม $out_dir ใน .latexmkrc)
%
%  ต่อจาก progress.tex (ครั้งที่ 1 · 7 ส.ค. 2569) ซึ่งยังเก็บไว้ไม่แก้ เพื่อให้
%  เทียบตัวเลขสองรอบกันได้ โครงเรื่องยึดตามสไลด์ Week 06 — PM: Project Progress
%  Report ของอาจารย์ เหมือนฉบับแรก
%
%  รอบนี้ต่างจากฉบับแรกตรงที่มี §10 "แผนที่ปรับปรุงแล้ว" ซึ่งเลื่อน Sprint ทั้งชุด
%  และภาคผนวก ข ที่เทียบแผนเดิมกับแผนใหม่ทีละรายการ ตัวเลขทุกตัวอ้างหลักฐานใน
%  repository ณ 10 ส.ค. 2569 โดย "ไม่นับ" โฟลเดอร์ spike-option-b/ ว่าเป็นผลงาน
%  ของผลิตภัณฑ์ เพราะเป็นงานทดลองสถาปัตยกรรม ไม่ใช่โค้ดที่ส่งมอบ
%
%  เขียนรายงานครั้งถัดไป: แก้ค่าใน "ตัวแปรของรอบรายงาน" ด้านล่าง แล้วไล่แก้
%  §3 (ตารางมิติงาน) · §4 (จุดบนกราฟ) · §5 (สถานะ Sprint) · §10 · ภาคผนวก ก
% =============================================================================
% =============================================================================
%  ulms-preamble.tex — preamble ที่ใช้ร่วมกันทุกเอกสารของโครงการ ULMs
%
%  ใช้โดย:  main.tex (ข้อเสนอโครงการ) · progress.tex (รายงานความก้าวหน้า)
%
%  เอกสารที่เรียกใช้ต้องทำสองอย่างหลัง % =============================================================================
%  ulms-preamble.tex — preamble ที่ใช้ร่วมกันทุกเอกสารของโครงการ ULMs
%
%  ใช้โดย:  main.tex (ข้อเสนอโครงการ) · progress.tex (รายงานความก้าวหน้า)
%
%  เอกสารที่เรียกใช้ต้องทำสองอย่างหลัง % =============================================================================
%  ulms-preamble.tex — preamble ที่ใช้ร่วมกันทุกเอกสารของโครงการ ULMs
%
%  ใช้โดย:  main.tex (ข้อเสนอโครงการ) · progress.tex (รายงานความก้าวหน้า)
%
%  เอกสารที่เรียกใช้ต้องทำสองอย่างหลัง \input{ulms-preamble}:
%    1) \hypersetup{pdftitle={...}}  — ชื่อเอกสารใน metadata ของ PDF
%    2) \begin{document}
%  (pdftitle ถูกย้ายออกจากไฟล์นี้ เพราะเป็นค่าเฉพาะของแต่ละเอกสาร)
%
%  แก้ที่นี่ที่เดียว = มีผลกับทุกเอกสาร ห้ามก๊อป preamble ไปวางซ้ำ
% =============================================================================
% =============================================================================
%  main.tex  —  เอกสารข้อเสนอโครงการ (Academic Project Proposal)
%  ระบบบริหารจัดการการยืม–คืนอุปกรณ์มหาวิทยาลัย (ULMs)
%
%  IMPORTANT: This document contains Thai text and MUST be compiled with
%  XeLaTeX (not pdflatex). Thai needs a Unicode font + ICU line-breaking.
%      Build:  latexmk main.tex        (uses .latexmkrc -> xelatex)
%      or:     xelatex main.tex        (run twice for the table of contents)
% =============================================================================

\documentclass[11pt, a4paper]{article}

% ---- Fonts & Thai support ----------------------------------------------------
% fontspec lets XeLaTeX use any system font. Laksaman is a complete serif face
% that covers BOTH Thai and Latin glyphs, so one \setmainfont handles the whole
% (Thai + English) document without font switching.
\usepackage{fontspec}

% Thai has no spaces between words, so TeX cannot find line-break points on its
% own. These two XeTeX primitives switch on ICU's dictionary-based Thai word
% segmentation, which is what makes Thai paragraphs wrap correctly.
\XeTeXlinebreaklocale "th"
\XeTeXlinebreakskip = 0pt plus 0.1pt

\setmainfont{Laksaman}

% ---- Page geometry & paragraph style ----------------------------------------
\usepackage[margin=2.5cm]{geometry}
\usepackage{parskip}          % block paragraphs (no indent, gap between) — common
                              % and more readable for Thai documents
\linespread{1.15}             % a little extra leading; Thai has tall tone marks

% ---- Colours (monochrome / grayscale palette) ------------------------------
\usepackage[table]{xcolor}    % 'table' option enables \rowcolor for table headers
% Monochrome (grayscale) palette. Names kept as-is so every reference still
% works; only the tones changed from the original green/blue accents.
\definecolor{kugreen}{HTML}{1A1A1A}   % near-black — section headings & title
\definecolor{kudark}{HTML}{333333}    % dark gray — subsection/subsubsection titles
\definecolor{headbg}{HTML}{EAEAEA}    % light gray background for table header rows
\definecolor{linkblue}{HTML}{404040}  % medium-dark gray — links / citations

% ---- Section heading styling -------------------------------------------------
% 'nobottomtitles*' stops a heading from being printed near the bottom of a
% page: if less than \bottomtitlespace of room is left, the heading (and its
% following text) moves to the next page. This is titlesec's built-in, robust
% replacement for the earlier \needspace hack — the latter injected vertical
% glue that clashed with longtable page-splitting ("Infinite glue shrinkage").
\usepackage[nobottomtitles*]{titlesec}
\titleformat{\section}
  {\color{kugreen}\Large\bfseries}{\thesection}{0.6em}{}
\titleformat{\subsection}
  {\color{kudark}\large\bfseries}{\thesubsection}{0.5em}{}
\titleformat{\subsubsection}
  {\color{kudark}\normalsize\bfseries}{\thesubsubsection}{0.5em}{}
\titlespacing*{\section}{0pt}{1.4\baselineskip}{0.5\baselineskip}

% Minimum room a heading needs at the bottom of a page; if only ~1–3 lines are
% left (less than this), the heading is bumped to the next page instead of being
% stranded above its table/text.
\renewcommand{\bottomtitlespace}{4\baselineskip}

% ---- Hierarchical indentation by heading level ------------------------------
% Each heading level shifts BOTH the heading and the body text under it further
% to the right: section = 0, subsection (x.y) = 1 step, subsubsection (x.y.z) =
% 2 steps. We do this by having the sectioning commands set \leftskip (a global
% left margin that carries through every following line until the next heading
% resets it). \par first, so the *previous* paragraph is not re-indented.
% \leftskip indents plain paragraphs; \@totalleftmargin makes list environments
% (itemize/enumerate) start from the same indent instead of the page margin.
\makeatletter
\newlength{\lvlindent}
\setlength{\lvlindent}{1.8em}          % one indent step
\let\ulmsSecIndent\section
\renewcommand{\section}{\par\global\leftskip=0pt\global\@totalleftmargin=0pt\relax\ulmsSecIndent}
\let\ulmsSubsecIndent\subsection
\renewcommand{\subsection}{\par\global\leftskip=\lvlindent\global\@totalleftmargin=\lvlindent\relax\ulmsSubsecIndent}
\let\ulmsSubsubsecIndent\subsubsection
\renewcommand{\subsubsection}{\par\global\leftskip=2\lvlindent\global\@totalleftmargin=2\lvlindent\relax\ulmsSubsubsecIndent}
\makeatother

% Indent the headings themselves to match (titlesec ignores \leftskip for the
% heading line, so its left margin is set here via the first \titlespacing arg).
\titlespacing*{\subsection}   {\lvlindent} {1.0\baselineskip}{0.4\baselineskip}
\titlespacing*{\subsubsection}{2\lvlindent}{0.8\baselineskip}{0.3\baselineskip}

% ---- Lists (tighter, custom bullets) ----------------------------------------
\usepackage{enumitem}
\setlist[itemize]{leftmargin=1.6em, itemsep=2pt, topsep=3pt}
\setlist[enumerate]{leftmargin=1.9em, itemsep=2pt, topsep=3pt}

% ---- Tables ------------------------------------------------------------------
\usepackage{array}
\usepackage{booktabs}         % \toprule / \midrule / \bottomrule
\usepackage{tabularx}         % auto-width columns that wrap long Thai text
\usepackage{longtable}        % tables that can break across pages
\usepackage{multirow}         % merged cells for the team-structure table

% Left-aligned wrapping column types:
%   L{width} — fixed-width ragged-right paragraph column
%   Y        — flexible ragged-right X column (for tabularx)
\newcolumntype{L}[1]{>{\raggedright\arraybackslash}p{#1}}
\newcolumntype{Y}{>{\raggedright\arraybackslash}X}
\renewcommand{\arraystretch}{1.3}

% Helpers for a styled header row inside tables.
%   \hrow  — starts a header row and shades it (must be FIRST in the row;
%            colortbl allows only one \rowcolor per row).
%   \thead — bold header-cell text.
\newcommand{\hrow}{\rowcolor{headbg}}
\newcommand{\thead}[1]{\textbf{#1}}

% \begin{keeptable} … \end{keeptable} — keeps a table and the bold label above
% it on the SAME page (used by the data dictionary in §5). A tabularx cannot
% split across pages, so LaTeX is free to break between the label and its
% table, stranding the label alone at the bottom of a page; wrapping the pair
% in a minipage makes them one unbreakable block.
%   * \leftskip is zeroed inside: the minipage box is already placed at the
%     current heading indent, so leaving it set would indent the content twice
%     and push the table past the right margin.
%   * \parskip is restored because minipage's \@parboxrestore zeroes it, which
%     would glue the table straight onto the label line.
% Tables inside must use \linewidth (= the minipage width), not \textwidth.
\newenvironment{keeptable}
  {\par\medskip\noindent
   \begin{minipage}{\dimexpr\textwidth-\leftskip\relax}%
   \leftskip=0pt\relax
   \setlength{\parskip}{0.5\baselineskip plus 2pt}}
  {\end{minipage}\par\medskip}

% ---- Table of contents styling ----------------------------------------------
% By default LaTeX gives dotted leaders to subsections but NOT to sections
% (they print bold with just a right-aligned page number). tocloft lets us turn
% on dots for the section level too, so every TOC line runs "….." to its page.
\usepackage{tocloft}
\renewcommand{\cftsecdotsep}{\cftdotsep}   % dotted leaders for section entries
\renewcommand{\cftsecleader}{\cftdotfill{\cftsecdotsep}}
\renewcommand{\cftsecpagefont}{\normalfont}  % keep page numbers un-bold

% ---- Figures -----------------------------------------------------------------
\usepackage{graphicx}
\graphicspath{{figures/}}

% Landscape (rotated) full-page floats — needed by the ER diagram in §5, which
% is far wider than it is tall and would be unreadable scaled to \textwidth.
\usepackage{rotating}

% ---- Flowcharts / diagrams (TikZ) -------------------------------------------
% Used for the borrow–return workflow diagram (§5). Libraries: geometric shapes
% for decision diamonds, arrows.meta for arrowheads, positioning for relative
% node placement, calc/fit/backgrounds for routing loops and grouping.
\usepackage{tikz}
\usetikzlibrary{shapes.geometric, arrows.meta, positioning, calc, fit, backgrounds}

% ---- Links & bookmarks -------------------------------------------------------
\usepackage{hyperref}
\hypersetup{
  colorlinks=true, linkcolor=black, citecolor=linkblue, urlcolor=linkblue,
  pdflang={th},
}

% Thai labels for auto-generated headings.
\renewcommand{\contentsname}{สารบัญ (Table of Contents)}
\renewcommand{\tablename}{ตาราง}
\renewcommand{\figurename}{รูปที่}

:
%    1) \hypersetup{pdftitle={...}}  — ชื่อเอกสารใน metadata ของ PDF
%    2) \begin{document}
%  (pdftitle ถูกย้ายออกจากไฟล์นี้ เพราะเป็นค่าเฉพาะของแต่ละเอกสาร)
%
%  แก้ที่นี่ที่เดียว = มีผลกับทุกเอกสาร ห้ามก๊อป preamble ไปวางซ้ำ
% =============================================================================
% =============================================================================
%  main.tex  —  เอกสารข้อเสนอโครงการ (Academic Project Proposal)
%  ระบบบริหารจัดการการยืม–คืนอุปกรณ์มหาวิทยาลัย (ULMs)
%
%  IMPORTANT: This document contains Thai text and MUST be compiled with
%  XeLaTeX (not pdflatex). Thai needs a Unicode font + ICU line-breaking.
%      Build:  latexmk main.tex        (uses .latexmkrc -> xelatex)
%      or:     xelatex main.tex        (run twice for the table of contents)
% =============================================================================

\documentclass[11pt, a4paper]{article}

% ---- Fonts & Thai support ----------------------------------------------------
% fontspec lets XeLaTeX use any system font. Laksaman is a complete serif face
% that covers BOTH Thai and Latin glyphs, so one \setmainfont handles the whole
% (Thai + English) document without font switching.
\usepackage{fontspec}

% Thai has no spaces between words, so TeX cannot find line-break points on its
% own. These two XeTeX primitives switch on ICU's dictionary-based Thai word
% segmentation, which is what makes Thai paragraphs wrap correctly.
\XeTeXlinebreaklocale "th"
\XeTeXlinebreakskip = 0pt plus 0.1pt

\setmainfont{Laksaman}

% ---- Page geometry & paragraph style ----------------------------------------
\usepackage[margin=2.5cm]{geometry}
\usepackage{parskip}          % block paragraphs (no indent, gap between) — common
                              % and more readable for Thai documents
\linespread{1.15}             % a little extra leading; Thai has tall tone marks

% ---- Colours (monochrome / grayscale palette) ------------------------------
\usepackage[table]{xcolor}    % 'table' option enables \rowcolor for table headers
% Monochrome (grayscale) palette. Names kept as-is so every reference still
% works; only the tones changed from the original green/blue accents.
\definecolor{kugreen}{HTML}{1A1A1A}   % near-black — section headings & title
\definecolor{kudark}{HTML}{333333}    % dark gray — subsection/subsubsection titles
\definecolor{headbg}{HTML}{EAEAEA}    % light gray background for table header rows
\definecolor{linkblue}{HTML}{404040}  % medium-dark gray — links / citations

% ---- Section heading styling -------------------------------------------------
% 'nobottomtitles*' stops a heading from being printed near the bottom of a
% page: if less than \bottomtitlespace of room is left, the heading (and its
% following text) moves to the next page. This is titlesec's built-in, robust
% replacement for the earlier \needspace hack — the latter injected vertical
% glue that clashed with longtable page-splitting ("Infinite glue shrinkage").
\usepackage[nobottomtitles*]{titlesec}
\titleformat{\section}
  {\color{kugreen}\Large\bfseries}{\thesection}{0.6em}{}
\titleformat{\subsection}
  {\color{kudark}\large\bfseries}{\thesubsection}{0.5em}{}
\titleformat{\subsubsection}
  {\color{kudark}\normalsize\bfseries}{\thesubsubsection}{0.5em}{}
\titlespacing*{\section}{0pt}{1.4\baselineskip}{0.5\baselineskip}

% Minimum room a heading needs at the bottom of a page; if only ~1–3 lines are
% left (less than this), the heading is bumped to the next page instead of being
% stranded above its table/text.
\renewcommand{\bottomtitlespace}{4\baselineskip}

% ---- Hierarchical indentation by heading level ------------------------------
% Each heading level shifts BOTH the heading and the body text under it further
% to the right: section = 0, subsection (x.y) = 1 step, subsubsection (x.y.z) =
% 2 steps. We do this by having the sectioning commands set \leftskip (a global
% left margin that carries through every following line until the next heading
% resets it). \par first, so the *previous* paragraph is not re-indented.
% \leftskip indents plain paragraphs; \@totalleftmargin makes list environments
% (itemize/enumerate) start from the same indent instead of the page margin.
\makeatletter
\newlength{\lvlindent}
\setlength{\lvlindent}{1.8em}          % one indent step
\let\ulmsSecIndent\section
\renewcommand{\section}{\par\global\leftskip=0pt\global\@totalleftmargin=0pt\relax\ulmsSecIndent}
\let\ulmsSubsecIndent\subsection
\renewcommand{\subsection}{\par\global\leftskip=\lvlindent\global\@totalleftmargin=\lvlindent\relax\ulmsSubsecIndent}
\let\ulmsSubsubsecIndent\subsubsection
\renewcommand{\subsubsection}{\par\global\leftskip=2\lvlindent\global\@totalleftmargin=2\lvlindent\relax\ulmsSubsubsecIndent}
\makeatother

% Indent the headings themselves to match (titlesec ignores \leftskip for the
% heading line, so its left margin is set here via the first \titlespacing arg).
\titlespacing*{\subsection}   {\lvlindent} {1.0\baselineskip}{0.4\baselineskip}
\titlespacing*{\subsubsection}{2\lvlindent}{0.8\baselineskip}{0.3\baselineskip}

% ---- Lists (tighter, custom bullets) ----------------------------------------
\usepackage{enumitem}
\setlist[itemize]{leftmargin=1.6em, itemsep=2pt, topsep=3pt}
\setlist[enumerate]{leftmargin=1.9em, itemsep=2pt, topsep=3pt}

% ---- Tables ------------------------------------------------------------------
\usepackage{array}
\usepackage{booktabs}         % \toprule / \midrule / \bottomrule
\usepackage{tabularx}         % auto-width columns that wrap long Thai text
\usepackage{longtable}        % tables that can break across pages
\usepackage{multirow}         % merged cells for the team-structure table

% Left-aligned wrapping column types:
%   L{width} — fixed-width ragged-right paragraph column
%   Y        — flexible ragged-right X column (for tabularx)
\newcolumntype{L}[1]{>{\raggedright\arraybackslash}p{#1}}
\newcolumntype{Y}{>{\raggedright\arraybackslash}X}
\renewcommand{\arraystretch}{1.3}

% Helpers for a styled header row inside tables.
%   \hrow  — starts a header row and shades it (must be FIRST in the row;
%            colortbl allows only one \rowcolor per row).
%   \thead — bold header-cell text.
\newcommand{\hrow}{\rowcolor{headbg}}
\newcommand{\thead}[1]{\textbf{#1}}

% \begin{keeptable} … \end{keeptable} — keeps a table and the bold label above
% it on the SAME page (used by the data dictionary in §5). A tabularx cannot
% split across pages, so LaTeX is free to break between the label and its
% table, stranding the label alone at the bottom of a page; wrapping the pair
% in a minipage makes them one unbreakable block.
%   * \leftskip is zeroed inside: the minipage box is already placed at the
%     current heading indent, so leaving it set would indent the content twice
%     and push the table past the right margin.
%   * \parskip is restored because minipage's \@parboxrestore zeroes it, which
%     would glue the table straight onto the label line.
% Tables inside must use \linewidth (= the minipage width), not \textwidth.
\newenvironment{keeptable}
  {\par\medskip\noindent
   \begin{minipage}{\dimexpr\textwidth-\leftskip\relax}%
   \leftskip=0pt\relax
   \setlength{\parskip}{0.5\baselineskip plus 2pt}}
  {\end{minipage}\par\medskip}

% ---- Table of contents styling ----------------------------------------------
% By default LaTeX gives dotted leaders to subsections but NOT to sections
% (they print bold with just a right-aligned page number). tocloft lets us turn
% on dots for the section level too, so every TOC line runs "….." to its page.
\usepackage{tocloft}
\renewcommand{\cftsecdotsep}{\cftdotsep}   % dotted leaders for section entries
\renewcommand{\cftsecleader}{\cftdotfill{\cftsecdotsep}}
\renewcommand{\cftsecpagefont}{\normalfont}  % keep page numbers un-bold

% ---- Figures -----------------------------------------------------------------
\usepackage{graphicx}
\graphicspath{{figures/}}

% Landscape (rotated) full-page floats — needed by the ER diagram in §5, which
% is far wider than it is tall and would be unreadable scaled to \textwidth.
\usepackage{rotating}

% ---- Flowcharts / diagrams (TikZ) -------------------------------------------
% Used for the borrow–return workflow diagram (§5). Libraries: geometric shapes
% for decision diamonds, arrows.meta for arrowheads, positioning for relative
% node placement, calc/fit/backgrounds for routing loops and grouping.
\usepackage{tikz}
\usetikzlibrary{shapes.geometric, arrows.meta, positioning, calc, fit, backgrounds}

% ---- Links & bookmarks -------------------------------------------------------
\usepackage{hyperref}
\hypersetup{
  colorlinks=true, linkcolor=black, citecolor=linkblue, urlcolor=linkblue,
  pdflang={th},
}

% Thai labels for auto-generated headings.
\renewcommand{\contentsname}{สารบัญ (Table of Contents)}
\renewcommand{\tablename}{ตาราง}
\renewcommand{\figurename}{รูปที่}

:
%    1) \hypersetup{pdftitle={...}}  — ชื่อเอกสารใน metadata ของ PDF
%    2) \begin{document}
%  (pdftitle ถูกย้ายออกจากไฟล์นี้ เพราะเป็นค่าเฉพาะของแต่ละเอกสาร)
%
%  แก้ที่นี่ที่เดียว = มีผลกับทุกเอกสาร ห้ามก๊อป preamble ไปวางซ้ำ
% =============================================================================
% =============================================================================
%  main.tex  —  เอกสารข้อเสนอโครงการ (Academic Project Proposal)
%  ระบบบริหารจัดการการยืม–คืนอุปกรณ์มหาวิทยาลัย (ULMs)
%
%  IMPORTANT: This document contains Thai text and MUST be compiled with
%  XeLaTeX (not pdflatex). Thai needs a Unicode font + ICU line-breaking.
%      Build:  latexmk main.tex        (uses .latexmkrc -> xelatex)
%      or:     xelatex main.tex        (run twice for the table of contents)
% =============================================================================

\documentclass[11pt, a4paper]{article}

% ---- Fonts & Thai support ----------------------------------------------------
% fontspec lets XeLaTeX use any system font. Laksaman is a complete serif face
% that covers BOTH Thai and Latin glyphs, so one \setmainfont handles the whole
% (Thai + English) document without font switching.
\usepackage{fontspec}

% Thai has no spaces between words, so TeX cannot find line-break points on its
% own. These two XeTeX primitives switch on ICU's dictionary-based Thai word
% segmentation, which is what makes Thai paragraphs wrap correctly.
\XeTeXlinebreaklocale "th"
\XeTeXlinebreakskip = 0pt plus 0.1pt

\setmainfont{Laksaman}

% ---- Page geometry & paragraph style ----------------------------------------
\usepackage[margin=2.5cm]{geometry}
\usepackage{parskip}          % block paragraphs (no indent, gap between) — common
                              % and more readable for Thai documents
\linespread{1.15}             % a little extra leading; Thai has tall tone marks

% ---- Colours (monochrome / grayscale palette) ------------------------------
\usepackage[table]{xcolor}    % 'table' option enables \rowcolor for table headers
% Monochrome (grayscale) palette. Names kept as-is so every reference still
% works; only the tones changed from the original green/blue accents.
\definecolor{kugreen}{HTML}{1A1A1A}   % near-black — section headings & title
\definecolor{kudark}{HTML}{333333}    % dark gray — subsection/subsubsection titles
\definecolor{headbg}{HTML}{EAEAEA}    % light gray background for table header rows
\definecolor{linkblue}{HTML}{404040}  % medium-dark gray — links / citations

% ---- Section heading styling -------------------------------------------------
% 'nobottomtitles*' stops a heading from being printed near the bottom of a
% page: if less than \bottomtitlespace of room is left, the heading (and its
% following text) moves to the next page. This is titlesec's built-in, robust
% replacement for the earlier \needspace hack — the latter injected vertical
% glue that clashed with longtable page-splitting ("Infinite glue shrinkage").
\usepackage[nobottomtitles*]{titlesec}
\titleformat{\section}
  {\color{kugreen}\Large\bfseries}{\thesection}{0.6em}{}
\titleformat{\subsection}
  {\color{kudark}\large\bfseries}{\thesubsection}{0.5em}{}
\titleformat{\subsubsection}
  {\color{kudark}\normalsize\bfseries}{\thesubsubsection}{0.5em}{}
\titlespacing*{\section}{0pt}{1.4\baselineskip}{0.5\baselineskip}

% Minimum room a heading needs at the bottom of a page; if only ~1–3 lines are
% left (less than this), the heading is bumped to the next page instead of being
% stranded above its table/text.
\renewcommand{\bottomtitlespace}{4\baselineskip}

% ---- Hierarchical indentation by heading level ------------------------------
% Each heading level shifts BOTH the heading and the body text under it further
% to the right: section = 0, subsection (x.y) = 1 step, subsubsection (x.y.z) =
% 2 steps. We do this by having the sectioning commands set \leftskip (a global
% left margin that carries through every following line until the next heading
% resets it). \par first, so the *previous* paragraph is not re-indented.
% \leftskip indents plain paragraphs; \@totalleftmargin makes list environments
% (itemize/enumerate) start from the same indent instead of the page margin.
\makeatletter
\newlength{\lvlindent}
\setlength{\lvlindent}{1.8em}          % one indent step
\let\ulmsSecIndent\section
\renewcommand{\section}{\par\global\leftskip=0pt\global\@totalleftmargin=0pt\relax\ulmsSecIndent}
\let\ulmsSubsecIndent\subsection
\renewcommand{\subsection}{\par\global\leftskip=\lvlindent\global\@totalleftmargin=\lvlindent\relax\ulmsSubsecIndent}
\let\ulmsSubsubsecIndent\subsubsection
\renewcommand{\subsubsection}{\par\global\leftskip=2\lvlindent\global\@totalleftmargin=2\lvlindent\relax\ulmsSubsubsecIndent}
\makeatother

% Indent the headings themselves to match (titlesec ignores \leftskip for the
% heading line, so its left margin is set here via the first \titlespacing arg).
\titlespacing*{\subsection}   {\lvlindent} {1.0\baselineskip}{0.4\baselineskip}
\titlespacing*{\subsubsection}{2\lvlindent}{0.8\baselineskip}{0.3\baselineskip}

% ---- Lists (tighter, custom bullets) ----------------------------------------
\usepackage{enumitem}
\setlist[itemize]{leftmargin=1.6em, itemsep=2pt, topsep=3pt}
\setlist[enumerate]{leftmargin=1.9em, itemsep=2pt, topsep=3pt}

% ---- Tables ------------------------------------------------------------------
\usepackage{array}
\usepackage{booktabs}         % \toprule / \midrule / \bottomrule
\usepackage{tabularx}         % auto-width columns that wrap long Thai text
\usepackage{longtable}        % tables that can break across pages
\usepackage{multirow}         % merged cells for the team-structure table

% Left-aligned wrapping column types:
%   L{width} — fixed-width ragged-right paragraph column
%   Y        — flexible ragged-right X column (for tabularx)
\newcolumntype{L}[1]{>{\raggedright\arraybackslash}p{#1}}
\newcolumntype{Y}{>{\raggedright\arraybackslash}X}
\renewcommand{\arraystretch}{1.3}

% Helpers for a styled header row inside tables.
%   \hrow  — starts a header row and shades it (must be FIRST in the row;
%            colortbl allows only one \rowcolor per row).
%   \thead — bold header-cell text.
\newcommand{\hrow}{\rowcolor{headbg}}
\newcommand{\thead}[1]{\textbf{#1}}

% \begin{keeptable} … \end{keeptable} — keeps a table and the bold label above
% it on the SAME page (used by the data dictionary in §5). A tabularx cannot
% split across pages, so LaTeX is free to break between the label and its
% table, stranding the label alone at the bottom of a page; wrapping the pair
% in a minipage makes them one unbreakable block.
%   * \leftskip is zeroed inside: the minipage box is already placed at the
%     current heading indent, so leaving it set would indent the content twice
%     and push the table past the right margin.
%   * \parskip is restored because minipage's \@parboxrestore zeroes it, which
%     would glue the table straight onto the label line.
% Tables inside must use \linewidth (= the minipage width), not \textwidth.
\newenvironment{keeptable}
  {\par\medskip\noindent
   \begin{minipage}{\dimexpr\textwidth-\leftskip\relax}%
   \leftskip=0pt\relax
   \setlength{\parskip}{0.5\baselineskip plus 2pt}}
  {\end{minipage}\par\medskip}

% ---- Table of contents styling ----------------------------------------------
% By default LaTeX gives dotted leaders to subsections but NOT to sections
% (they print bold with just a right-aligned page number). tocloft lets us turn
% on dots for the section level too, so every TOC line runs "….." to its page.
\usepackage{tocloft}
\renewcommand{\cftsecdotsep}{\cftdotsep}   % dotted leaders for section entries
\renewcommand{\cftsecleader}{\cftdotfill{\cftsecdotsep}}
\renewcommand{\cftsecpagefont}{\normalfont}  % keep page numbers un-bold

% ---- Figures -----------------------------------------------------------------
\usepackage{graphicx}
\graphicspath{{figures/}}

% Landscape (rotated) full-page floats — needed by the ER diagram in §5, which
% is far wider than it is tall and would be unreadable scaled to \textwidth.
\usepackage{rotating}

% ---- Flowcharts / diagrams (TikZ) -------------------------------------------
% Used for the borrow–return workflow diagram (§5). Libraries: geometric shapes
% for decision diamonds, arrows.meta for arrowheads, positioning for relative
% node placement, calc/fit/backgrounds for routing loops and grouping.
\usepackage{tikz}
\usetikzlibrary{shapes.geometric, arrows.meta, positioning, calc, fit, backgrounds}

% ---- Links & bookmarks -------------------------------------------------------
\usepackage{hyperref}
\hypersetup{
  colorlinks=true, linkcolor=black, citecolor=linkblue, urlcolor=linkblue,
  pdflang={th},
}

% Thai labels for auto-generated headings.
\renewcommand{\contentsname}{สารบัญ (Table of Contents)}
\renewcommand{\tablename}{ตาราง}
\renewcommand{\figurename}{รูปที่}



\hypersetup{pdftitle={รายงานความก้าวหน้าโครงการ ULMs ครั้งที่ 2}}

% ---- ตัวแปรของรอบรายงาน — แก้ที่นี่ที่เดียวเวลาออกรายงานรอบใหม่ ----------------
\newcommand{\reportno}{2}
\newcommand{\reportdate}{10 สิงหาคม 2569}
\newcommand{\reportperiod}{24 กรกฎาคม – 10 สิงหาคม 2569}
\newcommand{\plannedpct}{46\%}
\newcommand{\actualpct}{32\%}
\newcommand{\prevpct}{24\%}

\begin{document}

% =============================================================================
%  TITLE PAGE
% =============================================================================
\begin{titlepage}
\centering
\vspace*{2.5cm}

{\Large\bfseries\color{kugreen} รายงานความก้าวหน้าโครงการ ครั้งที่ \reportno\par}
\vspace{0.6cm}
{\LARGE\bfseries\color{kugreen} ระบบบริหารจัดการการยืม–คืนอุปกรณ์มหาวิทยาลัย\par}
\vspace{0.35cm}
{\Large\color{kudark} University Lending Management System (ULMs)\par}

\vspace{1.6cm}
\rule{0.72\textwidth}{0.8pt}
\vspace{0.9cm}

\begin{tabular}{@{}r@{\hspace{1.2em}}l@{}}
\textbf{รอบรายงาน}    & \reportperiod \\[0.35em]
\textbf{วันที่รายงาน}  & \reportdate \\[0.35em]
\textbf{สถานะ Sprint} & Sprint 1 (วันที่ 11 จาก 14) \\[0.35em]
\textbf{ความก้าวหน้า} & \actualpct\ เทียบแผน \plannedpct\ \ (ครั้งที่แล้ว \prevpct) \\[0.35em]
\textbf{สรุปสถานะ}    & \textbf{ช้ากว่าแผน · เสนอปรับแผน Sprint ทั้งชุด} \\
\end{tabular}

\vspace{0.9cm}
\rule{0.72\textwidth}{0.8pt}

\vfill
{\large\color{kudark} คณะทำงาน ``miro ปั่น 14 บาท''\par}
\vspace{0.3cm}
{\color{kudark} ภาควิชาวิศวกรรมคอมพิวเตอร์ คณะวิศวกรรมศาสตร์\par}
{\color{kudark} มหาวิทยาลัยเกษตรศาสตร์\par}
\vspace{1.2cm}
\end{titlepage}

\tableofcontents
\clearpage

% =============================================================================
\section{บทสรุปผู้บริหาร (Executive Summary)}

ณ วันที่ \reportdate\ โครงการอยู่ระหว่าง Sprint 1 เป็นวันที่ 11 จากทั้งหมด 14 วัน คิดเป็นระยะเวลาที่ใช้ไปแล้วประมาณ 2.4 สัปดาห์จากแผนทั้งหมด 10 สัปดาห์

\textbf{ความก้าวหน้าโดยรวมอยู่ที่ \actualpct\ เทียบกับแผนที่ควรอยู่ที่ \plannedpct\ นั่นคือช้ากว่าแผนประมาณ 14 จุด} ตัวเลขนี้ดีขึ้นจากรายงานครั้งที่ 1 ซึ่งอยู่ที่ \prevpct\ และช้ากว่าแผน 19 จุด แต่\textbf{ต้องอ่านตัวเลขนี้ด้วยความระมัดระวัง} เพราะเมื่อแยกดูเป็นรายมิติแล้วจะพบว่า \textbf{ส่วนต่างที่หายไป 5 จุดนั้นไม่ได้มาจากเส้นทางวิกฤต (Critical Path) แม้แต่จุดเดียว}

\subsection*{สิ่งที่เกิดขึ้นจริงในรอบสามวันที่ผ่านมา}

ในช่วงระหว่างรายงานครั้งที่ 1 กับฉบับนี้ คณะทำงานส่งมอบงานเพิ่มขึ้นจริงและเป็นงานที่มีคุณภาพ ได้แก่ \textbf{ระบบหน้าจอผู้ดูแลระบบ (Admin Console) จำนวน 6 หน้า} รวมประมาณ 2{,}000 บรรทัด \textbf{ระบบตรวจสอบอัตโนมัติ (CI)} จำนวน 3 งานตรวจ และ\textbf{ชุดเอกสารเตรียมสัญญา API} ที่ละเอียดมาก อย่างไรก็ตาม งานทั้งสามกลุ่มนี้มีลักษณะร่วมกันสองอย่างที่ต้องรายงานให้ชัด

ประการแรก \textbf{ไม่มีงานใดในสามกลุ่มนี้อยู่ในแผนของ Sprint 1} หน้าจอผู้ดูแลระบบตามแผนเดิมเป็นงานของ Sprint 3 ซึ่งถูกดึงขึ้นมาทำก่อนกำหนดราวสี่สัปดาห์ ส่วน CI เป็นงานตกค้างจาก Sprint 0 และเอกสารเตรียมสัญญาไม่เคยปรากฏในแผนฉบับใด

ประการที่สอง \textbf{ฝั่ง Back-End ไม่มีความคืบหน้าเพิ่มขึ้นเลยนับตั้งแต่วันที่ 5 สิงหาคม} การเปลี่ยนแปลงล่าสุดของโฟลเดอร์ \texttt{backend/} คือคอมมิต \texttt{ed18534} เมื่อวันที่ 5 สิงหาคม จนถึงวันที่ออกรายงานฉบับนี้ \textbf{ยังไม่มี Procedure ถูกเขียนขึ้นแม้แต่รายการเดียว ยังไม่มีระบบยืนยันตัวตน และยังไม่มีสคริปต์สร้างข้อมูลตั้งต้น} ฐานข้อมูลจึงยังว่างเปล่า และ \textbf{ยังไม่มีหน้าจอใดในระบบที่เรียกข้อมูลจากฐานข้อมูลจริงได้สักหน้าเดียว}

\subsection*{ข้อสรุปที่ต้องรายงานอย่างตรงไปตรงมา}

เป้าหมายส่งมอบของ Sprint 1 กำหนดไว้ว่าต้องยืมอุปกรณ์ข้ามภาควิชาได้ครบวงจร \textbf{เป้าหมายนี้ไม่มีทางสำเร็จภายในวันที่ 13 สิงหาคม 2569} เพราะงานฝั่ง Back-End ทั้งหมดที่จำเป็นยังไม่ได้เริ่ม คณะทำงานจึงเสนอ\textbf{ปรับแผน Sprint ทั้งชุด} ตามรายละเอียดในหัวข้อ \ref{sec:replan} โดยเลื่อนวงจรยืม–คืนครบรูปแบบไปเป็นเป้าหมายหลักของ Sprint 2 และเลื่อนงานที่เหลือถัดไปทีละหนึ่งชั้น

การเลื่อนหนึ่งชั้นนี้ทำให้\textbf{ไม่มี Sprint เหลือรองรับงานท้ายแผน} คณะทำงานจึงเสนอตัดขอบเขตสามรายการออกทันทีในรอบนี้ ไม่รอถึงจุดตัดสินใจเดิม

\subsection*{ข้อเสนอเร่งด่วนต่อคณะกรรมการ}

\begin{enumerate}
  \item \textbf{อนุมัติการปรับเกณฑ์ ``เสร็จ'' ของ Sprint 1} จาก ``ยืมข้ามภาคครบวงจร'' เหลือ ``มีหน้าจออย่างน้อยหนึ่งหน้าที่ดึงข้อมูลจากฐานข้อมูลจริงผ่าน tRPC ได้สำเร็จ'' เพื่อพิสูจน์ว่าทั้งสามชั้นของระบบต่อกันติดก่อน
  \item \textbf{ล็อกสัญญา tRPC ให้เสร็จภายในวันที่ 13 สิงหาคม 2569} งานเตรียมทั้งหมดพร้อมแล้ว ทั้งคู่มือ วาระประชุม 19 วาระ และโค้ดตัวอย่าง 28 ไฟล์ สิ่งที่ขาดคือการตัดสินใจในที่ประชุม ไม่ใช่เวลาเขียนโค้ด \textbf{ทุกงานของ Sprint 2 ขึ้นกับข้อนี้ข้อเดียว}
  \item \textbf{อนุมัติการตัดขอบเขต 3 รายการทันที} ได้แก่ การจองสถานที่ (Facility T3) การส่งออกรายงาน และการจองล่วงหน้าขั้นสูง เหตุผลและรายละเอียดอยู่ในหัวข้อ \ref{sec:scopechange}
  \item \textbf{เพิ่มกำลังคนฝั่ง Back-End} ซึ่งขณะนี้มีผู้ลงมือเขียนโค้ดจริงเพียงคนเดียว สวนทางกับสัดส่วนงานที่ฝั่ง Back-End รับผิดชอบสูงที่สุดในโครงการ ข้อเสนอนี้เคยยื่นไว้ในรายงานครั้งที่ 1 แล้วและยังไม่ได้รับการแก้ไข
\end{enumerate}

% =============================================================================
\section{สถานะโครงการโดยรวม (เงิน · งาน · เวลา)}

\begin{tabularx}{\dimexpr\textwidth-\leftskip\relax}{@{} L{2.2cm} L{3.1cm} L{3.1cm} Y @{}}
\toprule
\hrow \thead{ด้าน} & \thead{แผน} & \thead{ผลจริง} & \thead{การประเมิน} \\
\midrule
เวลา & ผ่านไป 2.4 จาก 10 สัปดาห์ (24\%) & ผ่านไป 2.4 จาก 10 สัปดาห์ (24\%) & เป็นไปตามปฏิทิน ยังไม่มีการเลื่อนวันสิ้นสุดโครงการ \\
งาน & \plannedpct & \actualpct & \textbf{ช้ากว่าแผน 14 จุด} ดีขึ้นจาก 19 จุดในรายงานครั้งที่แล้ว แต่ส่วนที่ดีขึ้นไม่ได้อยู่บนเส้นทางวิกฤต \\
เงิน & 5{,}000 บาท (วงเงินประเมินทั้งโครงการ) & 0 บาท (0\%) & ยังไม่มีรายจ่าย เพราะยังไม่ถึงขั้นเช่า Cloud/Hosting ตามที่ประเมินไว้ในข้อเสนอโครงการ \\
\bottomrule
\end{tabularx}

\vspace{0.8em}
\noindent ข้อสังเกตเดิมจากรายงานครั้งที่ 1 ยังใช้ได้อยู่ คือ \textbf{ค่าใช้จ่าย 0\% ไม่ใช่สัญญาณที่ดีในกรณีนี้} เพราะงบประมาณที่ตั้งไว้ผูกกับการนำระบบขึ้นใช้งานจริง การที่ยังไม่มีรายจ่ายเลยสะท้อนว่างานยังไม่เดินไปถึงขั้นนั้น ไม่ได้แปลว่าประหยัดงบได้ และเมื่อผ่านไปอีกหนึ่งรอบรายงานโดยตัวเลขยังเป็นศูนย์ ข้อสังเกตนี้ยิ่งมีน้ำหนักขึ้น

% =============================================================================
\section{ความก้าวหน้าแยกตามมิติของงาน}
\label{sec:dimension}

ตารางนี้ใช้วิธีคำนวณเดียวกับรายงานครั้งที่ 1 ทุกประการ เพื่อให้เทียบสองรอบกันได้ตรง ๆ ค่าน้ำหนักไม่เปลี่ยน และคอลัมน์สุดท้ายเพิ่มค่าของรอบที่แล้วไว้ให้เห็นทิศทาง ที่มาของทุกตัวเลขแจกแจงไว้ในภาคผนวก ก

\begin{tabularx}{\dimexpr\textwidth-\leftskip\relax}{@{} L{3.6cm} L{1.1cm} L{1.0cm} L{1.0cm} L{1.3cm} L{1.6cm} Y @{}}
\toprule
\hrow \thead{มิติงาน} & \thead{น้ำหนัก} & \thead{แผน} & \thead{จริง} & \thead{ผลต่าง} & \thead{รอบที่แล้ว} & \thead{หมายเหตุ} \\
\midrule
Back-End Design \newline (API + Database Model) & 10\% & 90\% & 55\% & $-35$ & 45\% & ดีขึ้นจากเอกสารเตรียมสัญญา แต่ยังไม่มีไฟล์สัญญาในรีโป \\
Front-End Design \newline (UI/UX + Business Process) & 10\% & 70\% & 72\% & $+2$ & 60\% & \textbf{เกินแผน} ออกแบบหน้าจอจริงแล้ว 7 จาก 26 หน้า \\
Back-End Dev & 30\% & 38\% & 12\% & $\mathbf{-26}$ & 12\% & \textbf{ไม่ขยับเลยตั้งแต่ 5 ส.ค.} ยังไม่มีตรรกะทางธุรกิจสักบรรทัด \\
Front-End Dev & 25\% & 44\% & 38\% & $-6$ & 20\% & เพิ่มขึ้นมากที่สุด จากหน้าจอผู้ดูแลระบบ 6 หน้า แต่ยังใช้ข้อมูลจำลองทั้งหมด \\
Test (Front-End + Back-End) & 15\% & 22\% & 0\% & $\mathbf{-22}$ & 0\% & \textbf{ยังไม่เริ่ม} ไม่มีทั้งแผนทดสอบและชุดทดสอบ \\
Documentation & 10\% & 45\% & 62\% & $+17$ & 45\% & \textbf{เร็วกว่าแผนมาก} เพิ่มเอกสารสัญญา 64 หน้าในรอบนี้ \\
\midrule
\hrow \thead{รวมถ่วงน้ำหนัก} & \thead{100\%} & \thead{\plannedpct} & \thead{\actualpct} & \thead{$\mathbf{-14}$} & \thead{\prevpct} & \\
\bottomrule
\end{tabularx}

\vspace{0.8em}
\noindent อ่านตารางนี้เทียบกับรายงานครั้งที่ 1 แล้วจะเห็นภาพที่ชัดกว่าตัวเลขรวม มิติที่ขยับขึ้นในรอบนี้มีสามมิติ คือ Front-End Dev (จาก 20\% เป็น 38\%) Documentation (จาก 45\% เป็น 62\%) และ Back-End Design (จาก 45\% เป็น 55\%) \textbf{ทั้งสามมิติมีน้ำหนักรวมกันเพียง 45\% และไม่มีมิติใดเป็นเงื่อนไขให้ Sprint ถัดไปเดินต่อได้}

ขณะเดียวกัน \textbf{มิติที่มีน้ำหนักสูงสุดคือ Back-End Dev (30\%) ยังคงอยู่ที่ 12\% เท่าเดิมทุกประการ} และมิติ Test (15\%) ยังคงเป็นศูนย์ สองมิตินี้รวมกันคิดเป็น 45\% ของน้ำหนักทั้งโครงการ และเป็นที่มาของช่องว่าง 14 จุดเกือบทั้งหมด

รูปแบบที่พบในรายงานครั้งที่ 1 จึงยังคงอยู่เหมือนเดิม คือ \textbf{งานที่ทำคนเดียวได้และไม่ต้องรอใครเดินหน้าได้ดี ส่วนงานที่ต้องอาศัยข้อตกลงร่วมกันก่อนจึงเริ่มได้ยังคงหยุดนิ่ง} ต่างกันเพียงว่ารอบนี้ ``งานที่ทำคนเดียวได้'' ขยายจากเอกสารและฐานข้อมูล มาถึงหน้าจอผู้ดูแลระบบด้วย

% =============================================================================
\clearpage
\section{กราฟความก้าวหน้าเทียบแผน (S-Curve)}

\begin{figure}[!ht]
\centering
\begin{tikzpicture}[x=1.05cm, y=0.055cm]
  % ---- กริดและแกน ----
  \foreach \y in {0,20,40,60,80,100}
    \draw[gray!25] (0,\y) -- (10.4,\y);
  \draw[->, thick] (0,0) -- (10.8,0) node[right, font=\small] {สัปดาห์};
  \draw[->, thick] (0,0) -- (0,112) node[above, font=\small] {\% ความก้าวหน้า};
  \foreach \y in {0,20,40,60,80,100}
    \draw (0,\y) -- (-0.12,\y) node[left, font=\scriptsize] {\y};
  \foreach \x/\lbl in {0/{24 ก.ค.}, 2/{7 ส.ค.}, 4/{21 ส.ค.}, 6/{4 ก.ย.}, 8/{18 ก.ย.}, 10/{1 ต.ค.}}
    \draw (\x,0) -- (\x,-2) node[below, font=\scriptsize] {\lbl};

  % ---- เส้นแผน (Planned) ----
  \draw[thick, kudark, dashed]
    plot coordinates {(0,0) (1,25) (3,55) (5,75) (7,90) (9,97) (10,100)};
  \foreach \p in {(0,0), (1,25), (3,55), (5,75), (7,90), (9,97), (10,100)}
    \fill[kudark] \p circle (1.6pt);

  % ---- เส้นผลจริง (Actual) — จุดสุดท้ายคือรายงานฉบับนี้ ----
  \draw[very thick, kugreen]
    plot coordinates {(0,0) (1,14) (2,24) (2.4,32)};
  \foreach \p in {(0,0), (1,14), (2,24), (2.4,32)}
    \fill[kugreen] \p circle (2.2pt);

  % ---- ป้ายกำกับจุดของรายงานสองรอบ ----
  % วางไว้ใต้เส้นผลจริงทั้งคู่ เพื่อไม่ให้ทับเส้นแผนที่พาดอยู่ด้านบน
  \node[below right, font=\scriptsize, gray!70!black] at (2.05,23)
    {ครั้งที่ 1 (7 ส.ค.) 24\%};
  \node[below right, font=\scriptsize, kugreen] at (2.55,32)
    {ครั้งที่ 2 (10 ส.ค.) 32\%};

  % ---- ช่องว่างระหว่างแผนกับผลจริง ณ วันที่รายงาน ----
  \draw[<->, red!70!black, thick] (2.4,32) -- (2.4,46);
  \node[right, font=\scriptsize, red!70!black] at (2.5,41) {ช้ากว่าแผน 14 จุด};
  \fill[red!70!black] (2.4,46) circle (1.6pt);

  % ---- คำอธิบายสัญลักษณ์ ----
  \draw[thick, kudark, dashed] (5.6,18) -- (6.4,18);
  \node[right, font=\scriptsize] at (6.5,18) {แผน (Planned)};
  \draw[very thick, kugreen] (5.6,8) -- (6.4,8);
  \node[right, font=\scriptsize] at (6.5,8) {ผลจริง (Actual)};
\end{tikzpicture}
\caption{กราฟ S-Curve เปรียบเทียบความก้าวหน้าตามแผนกับผลการดำเนินงานจริง
  เส้นผลจริงเพิ่มจุดที่สัปดาห์ที่ 2.4 ซึ่งเป็นวันที่ออกรายงานฉบับนี้
  ช่องว่างแคบลงจาก 19 จุดเหลือ 14 จุด
  ค่าความก้าวหน้าคำนวณแบบถ่วงน้ำหนักตามตารางในหัวข้อ \ref{sec:dimension}}
\label{fig:scurve}
\end{figure}

\noindent รูปที่ \ref{fig:scurve} แสดงให้เห็นว่าเส้นผลจริงชันขึ้นในช่วงสามวันหลัง ซึ่งเป็นสัญญาณที่ดีในตัวเอง แต่\textbf{ความชันนี้มาจากมิติที่มีน้ำหนักรวมกันไม่ถึงครึ่ง} หากมิติ Back-End Dev ยังคงนิ่งอยู่ที่ 12\% ต่อไป เส้นผลจริงจะชนเพดานที่ประมาณ 55\% ไม่ว่าจะทำงานฝั่งอื่นเพิ่มอีกเท่าไรก็ตาม เพราะน้ำหนักที่เหลือทั้งหมดผูกอยู่กับงานที่ต้องมี Back-End ก่อน

% =============================================================================
\clearpage
\section{แผนและผลการดำเนินงานที่ผ่านมา}

\subsection{Sprint 0 ``วางฐาน'' (24 – 30 กรกฎาคม 2569)}

รายงานครั้งที่ 1 ประเมิน Sprint 0 ด้วยเกณฑ์ ``เสร็จ / ไม่เสร็จ'' เพียงสองค่า ซึ่งทำให้งานที่เสร็จไปแล้วส่วนใหญ่ถูกนับเป็นศูนย์ รอบนี้จึงประเมินใหม่ด้วยเกณฑ์ที่ให้คะแนนบางส่วนได้ ผลคือ Sprint 0 อยู่ที่ประมาณ 77\% ไม่ใช่ 27\% อย่างที่ตารางเดิมแสดง

\begin{tabularx}{\dimexpr\textwidth-\leftskip\relax}{@{} L{1.1cm} Y L{2.6cm} @{}}
\toprule
\hrow \thead{ส่วนงาน} & \thead{งานตามแผน} & \thead{สถานะจริง} \\
\midrule
SA & ปิดขอบเขตของโครงการ & \textbf{เสร็จ} \\
BE & Prisma schema แกนกลาง 33 ตาราง พร้อม Migration ขึ้นฐานข้อมูลจริง & \textbf{85\%} ขาดตารางผูกชิ้นอุปกรณ์และสถานะอนุมัติบางส่วน \\
BE & สัญญา tRPC ของวงจรหลัก แบบอิงสมาชิกตามหน่วยงาน & \textbf{25\%} มีเอกสารเตรียมครบ แต่ไม่มีไฟล์สัญญาในรีโป \\
BE & ติดตั้ง NestJS + tRPC + Prisma + PostgreSQL พร้อม Migration & \textbf{เสร็จ} \\
BE & ข้อมูลตั้งต้น 2 หน่วยงาน & \textbf{เลื่อนไป Sprint 2} \\
FE & ติดตั้ง Vite + React + TypeScript + Router และโครงยืนยันตัวตน & \textbf{เสร็จ} \\
FE & ระบบตรวจสอบอัตโนมัติ (CI) & \textbf{80\%} เขียนเสร็จแล้วบนสาขา \texttt{ci-setup} แต่ยังไม่รวมเข้าสายหลัก \\
FE & ระบบออกแบบและ Wireframe หน้าเข้าสู่ระบบกับเปลือกระบบ & \textbf{เสร็จ} \\
FE & ชั้นข้อมูลจำลองที่สร้างจากสัญญา API & \textbf{55\%} มีข้อมูลจำลองใช้งานจริง แต่เขียนด้วยมือ ไม่ได้สร้างจากสัญญา \\
QA & แผนการทดสอบและเกณฑ์การยอมรับ & \textbf{เลื่อนไป Sprint 2} \\
PM & ตั้งกระดานงาน พิธีการประจำ Sprint และเกณฑ์ ``เสร็จ'' & \textbf{เสร็จ} \\
\bottomrule
\end{tabularx}

\vspace{0.8em}
\noindent\textbf{เกณฑ์ ``เสร็จ'' ของ Sprint 0} ยังไม่ผ่าน เพราะขาดสองข้อเดิม คือสัญญา API ที่ยังไม่ล็อก และข้อมูลตั้งต้นที่ยังไม่มี ส่วนข้อที่ว่า Repository ต้องรันได้ทั้งสองฝั่งนั้นผ่านแล้ว และได้รับการยืนยันซ้ำจากระบบ CI ที่เขียนขึ้นในรอบนี้ ซึ่งตรวจสอบการติดตั้ง การตรวจชนิดข้อมูล และการ Build ของทั้งสองฝั่ง

\subsection{Sprint 1 (31 กรกฎาคม – 13 สิงหาคม 2569)}

\noindent\textbf{เป้าหมายตามแผนเดิม} ทำให้เกิดเส้นทางการทำงานที่ทะลุทุกชั้นของระบบ (Vertical Slice) โดยยืมอุปกรณ์ข้ามหน่วยงานได้ในเส้นทางปกติ

\begin{tabularx}{\dimexpr\textwidth-\leftskip\relax}{@{} L{1.1cm} Y L{2.6cm} @{}}
\toprule
\hrow \thead{ส่วนงาน} & \thead{งานตามแผน} & \thead{สถานะจริง} \\
\midrule
BE & ยืนยันตัวตนและควบคุมสิทธิ์แบบอิงสมาชิกต่อหน่วยงาน & \textbf{ยังไม่เริ่ม} \\
BE & รายการอุปกรณ์ ค้นหา จำนวนคงเหลือ และการแสดงหน่วยงานเจ้าของ & \textbf{ยังไม่เริ่ม} \\
BE & สร้างคำขอ อนุมัติอัตโนมัติ ผูกชิ้น ส่งมอบ และรับคืน & \textbf{ยังไม่เริ่ม} \\
FE & หน้าเข้าสู่ระบบและการคุ้มกันเส้นทาง & \textbf{เสร็จ} (ยังใช้การยืนยันตัวตนจำลอง) \\
FE & หน้าเรียกดูอุปกรณ์ รายละเอียด และจำนวนคงเหลือ & \textbf{ยังไม่เริ่ม} \\
FE & หน้าส่งคำขอและติดตามสถานะของผู้ยืม & \textbf{ยังไม่เริ่ม} \\
FE & หน้าคิวงานของเจ้าหน้าที่ & \textbf{ยังไม่เริ่ม} \\
QA & ทดสอบเส้นทางปกติและกรณียืมข้ามภาควิชา & \textbf{ยังไม่เริ่ม} \\
\midrule
\hrow \multicolumn{3}{@{}l@{}}{\thead{งานที่ทำจริงในช่วง Sprint 1 แต่ไม่ได้อยู่ในแผน}} \\
FE & หน้าจอผู้ดูแลระบบ 6 หน้า (แผนเดิมอยู่ใน Sprint 3) & \textbf{เสร็จ} (ใช้ข้อมูลจำลอง) \\
FE & ระบบตรวจสอบอัตโนมัติ CI (งานตกค้างจาก Sprint 0) & \textbf{80\%} \\
SA & เอกสารเตรียมสัญญา tRPC 64 หน้า และโค้ดตัวอย่าง 28 ไฟล์ & \textbf{เสร็จ} \\
\bottomrule
\end{tabularx}

\vspace{0.8em}
\noindent ผ่านไปแล้ว 11 วันจาก 14 วันของ Sprint 1 \textbf{ยังไม่มีงานใดในรายการตามแผนที่เสร็จสมบูรณ์ นอกจากหน้าเข้าสู่ระบบซึ่งทำงานอยู่บนการยืนยันตัวตนจำลอง} ขณะที่งานที่เสร็จจริงในช่วงเวลาเดียวกันนั้นเป็นงานนอกแผนทั้งสามรายการ

ข้อสังเกตที่ต้องบันทึกไว้คือ \textbf{งานนอกแผนทั้งสามรายการไม่ใช่งานเสียเปล่า} หน้าจอผู้ดูแลระบบเป็นงานที่ต้องทำอยู่แล้ว การทำก่อนกำหนดจึงเป็นการดึงงานอนาคตมาเก็บ ระบบ CI เป็นเครื่องมือที่จะป้องกันไม่ให้โค้ดเสียหลุดเข้าสายหลักตลอดโครงการ และเอกสารเตรียมสัญญาลดเวลาที่ต้องใช้ในการตัดสินใจในที่ประชุมลงมาก \textbf{ปัญหาจึงไม่ใช่คุณค่าของงาน แต่เป็นลำดับก่อนหลัง} คือทีมเลือกทำงานที่ทำได้ทันทีโดยไม่ต้องรอใคร แทนที่จะทำงานที่ปลดล็อกให้คนอื่นเดินต่อได้

% =============================================================================
\clearpage
\section{ผลงานที่ส่งมอบแล้ว}

หัวข้อนี้รายงานเฉพาะผลงานที่ตรวจสอบได้จริงใน Repository ไม่รวมกิจกรรมระหว่างดำเนินงาน และ\textbf{ไม่นับโฟลเดอร์ \texttt{spike-option-b/}} ซึ่งเป็นงานทดลองสถาปัตยกรรม ไม่ใช่โค้ดที่ส่งมอบเป็นตัวผลิตภัณฑ์

\subsection{ฝั่ง Front-End}

\begin{itemize}
  \item \textbf{โครงสร้างพื้นฐานพร้อมใช้งาน} Vite + React 18 + TypeScript + React Router 6 พร้อมเส้นทางของหน้าจอครบ 26 เส้นทาง และชั้นคุ้มกันสิทธิ์ตามบทบาท 4 ชั้นซ้อนอยู่บนเส้นทางเหล่านั้น
  \item \textbf{ระบบออกแบบและเปลือกระบบ} Tailwind CSS พร้อมโทเคนสี แถบนำทางด้านข้าง แถบด้านบน เมนูที่เปลี่ยนตามบทบาททั้งสี่บทบาท และชุดส่วนประกอบพื้นฐาน 19 ตัว ซึ่งรวมตารางข้อมูล กราฟ กล่องโต้ตอบ และแผงเลื่อน
  \item \textbf{รองรับสองภาษาและสองธีม} ไทย–อังกฤษผ่าน i18next พร้อมโหมดสว่างและมืด
  \item \textbf{หน้าเข้าสู่ระบบที่ใช้งานได้จริง} รองรับสองช่องทาง คือผ่านอีเมล \texttt{@ku.ac.th} และบัญชีภายใน พร้อมการคุ้มกันเส้นทางตามบทบาท
  \item \textbf{หน้าจอผู้ดูแลระบบ 6 หน้า} (\textbf{ใหม่ในรอบนี้}) ได้แก่ แผงสรุป ผู้ใช้ สถานะระบบ บันทึกการตรวจสอบ ตั้งค่า และรายงาน รวมประมาณ 2{,}000 บรรทัด โดยหน้าผู้ใช้เพียงหน้าเดียวมี 511 บรรทัด
\end{itemize}

\noindent\textbf{ข้อจำกัดที่ต้องระบุให้ชัด} จากทั้งหมด 26 หน้า มี \textbf{7 หน้าที่มีเนื้อหาจริง และอีก 19 หน้ายังเป็นหน้าเปล่า} (Placeholder ที่มีเพียง 6 บรรทัดต่อหน้า) ซึ่งครอบคลุมหน้าของผู้ยืมทั้งหมด 7 หน้า หน้าของเจ้าหน้าที่ทั้งหมด 7 หน้า และหน้าของผู้อนุมัติทั้งสองหน้า

ข้อจำกัดที่สำคัญกว่าคือ \textbf{ทั้ง 7 หน้าที่มีเนื้อหาจริงนั้นอ่านข้อมูลจากไฟล์จำลองในโค้ดทั้งหมด ไม่มีหน้าใดเรียก Back-End เลย} ไฟล์ \texttt{lib/api-client.ts} ซึ่งเขียนไว้สำหรับเรียก Back-End นั้นมีอยู่จริงและเขียนเสร็จแล้ว แต่ตรวจสอบทั้งโฟลเดอร์แล้ว\textbf{ไม่พบว่ามีไฟล์ใดเรียกใช้งานมันเลยแม้แต่ไฟล์เดียว}

\subsection{ฝั่ง Back-End}

\begin{itemize}
  \item \textbf{โครงสร้างพื้นฐานพร้อมใช้งาน} NestJS พร้อม Prisma, \texttt{nestjs-trpc}, ตัวเชื่อม PostgreSQL และ \texttt{ConfigModule} ประกอบเข้าด้วยกันใน \texttt{app.module.ts} เรียบร้อย
  \item \textbf{โครงสร้างฐานข้อมูลครบถ้วนและใช้งานได้จริง} \texttt{schema.prisma} นิยามตารางครบทั้ง 33 ตาราง และมี Migration ที่สร้างตารางทั้งหมดบนฐานข้อมูลจริงแล้วเมื่อวันที่ 5 สิงหาคม 2569
\end{itemize}

\noindent\textbf{ข้อจำกัดที่ต้องระบุให้ชัด} รายการข้างต้น\textbf{เหมือนกับรายงานครั้งที่ 1 ทุกตัวอักษร} เพราะโฟลเดอร์ \texttt{backend/} ไม่มีการเปลี่ยนแปลงใดเลยนับตั้งแต่วันที่ 5 สิงหาคม 2569 \textbf{ยังไม่มี Procedure ใดถูกเขียนขึ้นแม้แต่รายการเดียว ยังไม่มีระบบยืนยันตัวตน ยังไม่มีตรรกะทางธุรกิจในโฟลเดอร์ \texttt{src/} และยังไม่มีสคริปต์สร้างข้อมูลตั้งต้น} ฐานข้อมูลที่สร้างขึ้นจึงยังว่างเปล่า

\subsection{ฝั่งเครื่องมือและกระบวนการ (ใหม่ในรอบนี้)}

\begin{itemize}
  \item \textbf{ระบบตรวจสอบอัตโนมัติ (CI)} ไฟล์ \texttt{.github/workflows/ci.yml} กำหนดงานตรวจ 3 งาน คือ ตรวจว่าทั้งสองฝั่งใช้ไลบรารีร่วมเวอร์ชันเดียวกัน ตรวจชนิดข้อมูลและ Build ฝั่ง Back-End และตรวจชนิดข้อมูลและ Build ฝั่ง Front-End พร้อมสคริปต์ช่วยตรวจ \texttt{scripts/check-versions.mjs}
  \item \textbf{ยังไม่ถูกรวมเข้าสายหลัก} ไฟล์ทั้งสองยังอยู่บนสาขา \texttt{ci-setup} และยังไม่ได้คอมมิต จึงยังไม่มีผลบังคับใช้กับคำขอรวมโค้ดใด ๆ ในขณะนี้
\end{itemize}

\subsection{ฝั่งเอกสารและการออกแบบ}

\begin{itemize}
  \item \textbf{ข้อเสนอโครงการฉบับสมบูรณ์ 42 หน้า} และ\textbf{การออกแบบฐานข้อมูลรอบที่ 4} พร้อมพจนานุกรมข้อมูลครบ 33 ตาราง (รายงานไว้แล้วในครั้งที่ 1)
  \item \textbf{คู่มือสัญญา tRPC 31 หน้า} (\textbf{ใหม่}) อธิบายหลักการและระบุความเสี่ยง 8 ข้อที่พบใน Repository นี้โดยเฉพาะ
  \item \textbf{วาระประชุมเรื่องสัญญา 33 หน้า 19 วาระ} (\textbf{ใหม่}) แต่ละวาระระบุทางเลือกและผลกระทบไว้ครบ เพื่อให้ที่ประชุมตัดสินใจได้โดยไม่ต้องค้นเพิ่ม
  \item \textbf{โค้ดตัวอย่างผลลัพธ์ของสัญญา 28 ไฟล์} (\textbf{ใหม่}) แสดงว่าโค้ดจะมีหน้าตาอย่างไรเมื่อวาระถูกตัดสินแล้ว พร้อมตารางแมประหว่างวาระกับไฟล์
  \item \textbf{สารบัญไฟล์ของโครงการ} (\texttt{docs/INDEX.md}) ที่ระบุว่าเอกสารฉบับล่าสุดของแต่ละเรื่องอยู่ที่ไหน ซึ่งจำเป็นเพราะโฟลเดอร์ \texttt{report/} มีเอกสารปนกันหลายรุ่น
\end{itemize}

% =============================================================================
\section{การเปลี่ยนแปลงขอบเขตโครงการ}
\label{sec:scopechange}

\subsection{การตัดขอบเขตที่ขอมติในรอบนี้}

การเลื่อนงานทั้งชุดไปหนึ่ง Sprint ทำให้ไม่มี Sprint เหลือรองรับงานท้ายแผน คณะทำงานจึงเสนอตัดสามรายการต่อไปนี้ออกทันที ไม่รอถึงจุดตัดสินใจวันที่ 20 สิงหาคมตามที่เคยเสนอไว้ เหตุผลคือการตัดตอนนี้ทำให้วางแผน Sprint 3 และ 4 ได้ทันที ส่วนการรอไปตัดทีหลังจะทำให้ทีมต้องเผื่อเวลาให้งานที่สุดท้ายไม่ได้ทำอยู่ดี

\begin{tabularx}{\dimexpr\textwidth-\leftskip\relax}{@{} L{4.0cm} Y L{3.0cm} @{}}
\toprule
\hrow \thead{รายการที่ขอตัด} & \thead{เหตุผล} & \thead{ผลกระทบต่อการนำเสนอ} \\
\midrule
การจองสถานที่ (Facility T3) & ถูกระบุไว้ตั้งแต่ต้นว่าเป็นงานที่ตัดได้ และไม่มีงานอื่นขึ้นกับมัน & ไม่มี ยังคงสาธิตการยืมอุปกรณ์และการจองล่วงหน้าซึ่งเป็นจุดเด่นหลักได้ครบ \\
การส่งออกรายงาน (Export) & หน้าจอรายงานของผู้ดูแลระบบทำเสร็จแล้วในรอบนี้ ใช้แสดงในการนำเสนอแทนได้ & ไม่มี \\
การจองล่วงหน้าขั้นสูง & ได้แก่ การผูกชิ้นอุปกรณ์อัตโนมัติเมื่อใกล้วันรับ การจัดคิวเมื่อจองไม่สำเร็จ และเพดานการต่ออายุที่เทียบกับวันจอง & \textbf{ปานกลาง} ยังจองล่วงหน้าได้ แต่เจ้าหน้าที่ต้องผูกชิ้นอุปกรณ์ด้วยมือ \\
\bottomrule
\end{tabularx}

\subsection{การเปลี่ยนแปลงอื่นที่ต้องบันทึก}

\begin{tabularx}{\dimexpr\textwidth-\leftskip\relax}{@{} L{4.0cm} Y L{2.4cm} @{}}
\toprule
\hrow \thead{เรื่อง} & \thead{การเปลี่ยนแปลง} & \thead{สถานะ} \\
\midrule
หน้าจอผู้ดูแลระบบ & แผนเดิมกำหนดไว้ใน Sprint 3 แต่ทีมทำเสร็จแล้วในช่วง Sprint 1 & \textbf{ดึงขึ้นมาแล้ว} ต้องเปลี่ยนจากข้อมูลจำลองเป็นข้อมูลจริงใน Sprint 2 \\
ตารางผูกชิ้นอุปกรณ์ (ALLOCATION) & รายงานไว้แล้วในครั้งที่ 1 ว่าขาด และ\textbf{ยังคงขาดอยู่} ในสคีมา 33 ตารางที่ Migrate ไปแล้ว & \textbf{ต้องแก้ใน Sprint 2} \\
สถานะ ``อนุมัติบางส่วน'' & รายงานไว้แล้วในครั้งที่ 1 ว่าขาด และ\textbf{ยังคงขาดอยู่} & \textbf{ต้องแก้ใน Sprint 2} \\
คีย์ไม่ซ้ำบนอีเมลและรหัสผู้ใช้ & \texttt{AccountInfo.Email} และ \texttt{UserID} ไม่มีข้อกำหนดห้ามซ้ำ ฐานข้อมูลจึงยอมให้สมัครอีเมลเดียวกันได้หลายบัญชี ซึ่งสำหรับระบบเข้าสู่ระบบถือเป็นช่องโหว่ & \textbf{ต้องแก้ใน Sprint 2} \\
ฐานข้อมูล MongoDB & ยังไม่มีข้อสรุป เป็นข้อค้างจากรายงานครั้งที่ 1 & รอคำแนะนำจากอาจารย์ \\
\bottomrule
\end{tabularx}

\vspace{0.8em}
\noindent ข้อสังเกตคือ \textbf{สามรายการแรกในตารางที่สองนี้ถูกรายงานไว้ตั้งแต่ครั้งที่ 1 และยังไม่มีรายการใดได้รับการแก้ไข} ต้นทุนการแก้จะยังต่ำอยู่ตราบใดที่ฐานข้อมูลยังว่างเปล่า และจะสูงขึ้นทันทีที่มีข้อมูลจริงและมีโค้ดเรียกใช้ตารางเหล่านั้นแล้ว

% =============================================================================
\clearpage
\section{ปัญหาอุปสรรคและแนวทางแก้ไข}

\begin{tabularx}{\dimexpr\textwidth-\leftskip\relax}{@{} L{3.2cm} Y L{2.5cm} @{}}
\toprule
\hrow \thead{ปัญหา} & \thead{ผลกระทบและแนวทางแก้ไข} & \thead{ผู้รับผิดชอบ / กำหนด} \\
\midrule
สัญญา tRPC ยังไม่ถูกล็อก \newline \textbf{(ค้างจากครั้งที่ 1)} & \textbf{ผลกระทบ} ยังเป็นต้นเหตุของงานที่ตกแผนแทบทั้งหมดเหมือนเดิม \newline \textbf{สิ่งที่เปลี่ยนไป} รอบนี้งานเตรียมเสร็จหมดแล้ว ทั้งคู่มือ วาระประชุม 19 วาระ และโค้ดตัวอย่าง 28 ไฟล์ \textbf{สิ่งที่ขาดจึงไม่ใช่เวลาเขียน แต่เป็นการนัดประชุมและลงมติ} \newline \textbf{แนวทาง} จัดประชุมลงมติวาระที่จำเป็นต่อวงจรหลักก่อน แล้วคอมมิตไฟล์สัญญาเข้ารีโปในวันเดียวกัน & API + SA \newline ภายใน 13 ส.ค. \\
ฝั่ง Back-End ไม่ขยับมา 5 วัน \newline \textbf{(รุนแรงขึ้น)} & \textbf{ผลกระทบ} เป้าหมายของ Sprint 1 ไม่สำเร็จแน่นอน และ Sprint 2 จะแบกงานสอง Sprint รวมกัน \newline \textbf{แนวทาง} แยกงานที่ทำได้ทันทีโดยไม่ต้องรอสัญญาออกมาทำก่อน ได้แก่ สคริปต์ข้อมูลตั้งต้น การเพิ่มตารางที่ขาด และคีย์ไม่ซ้ำ ทั้งสามอย่างไม่ขึ้นกับสัญญาเลย & BE-lead + BE-dev \newline เริ่มทันที \\
ไม่มีข้อมูลตั้งต้นในฐานข้อมูล \newline \textbf{(ค้างจากครั้งที่ 1)} & \textbf{ผลกระทบ} ต่อให้เขียน Procedure เสร็จก็ทดสอบไม่ได้ และเป็นงานที่กำหนดไว้ว่าต้องเสร็จตั้งแต่ 10 ส.ค. \newline \textbf{แนวทาง} เป็นงานแรกของ Sprint 2 ไม่ต้องรออะไรทั้งสิ้น & BE-dev \newline ภายใน 15 ส.ค. \\
งานทดสอบยังไม่เริ่มเลย \newline \textbf{(ค้างจากครั้งที่ 1)} & \textbf{ผลกระทบ} มิติ Test มีน้ำหนัก 15\% และยังเป็นศูนย์มาสองรอบรายงานติดกัน ยิ่งเริ่มช้า ค่าใช้จ่ายในการแก้ข้อบกพร่องยิ่งสูง \newline \textbf{แนวทาง} เกณฑ์การยอมรับเขียนได้ทันทีโดยไม่ต้องรอโค้ด ให้เริ่มจากวงจรหลักก่อน & QA-lead \newline ภายใน 17 ส.ค. \\
CI เขียนเสร็จแต่ยังไม่บังคับใช้ \newline \textbf{(ใหม่)} & \textbf{ผลกระทบ} เครื่องมือที่ป้องกันโค้ดเสียมีอยู่แล้วแต่ไม่ทำงาน เพราะยังไม่ถูกรวมเข้าสายหลัก \newline \textbf{แนวทาง} เปิดคำขอรวมโค้ดทันที ซึ่งจะทำให้ CI ทดสอบตัวเองในคำขอนั้นไปด้วย & FE-lead \newline ภายใน 14 ส.ค. \\
ลำดับความสำคัญของงาน \newline \textbf{(ใหม่)} & \textbf{ผลกระทบ} รอบนี้ทีมส่งมอบงานได้จริงสามรายการ แต่ทั้งสามอยู่นอกแผนและไม่ปลดล็อกใครต่อ ขณะที่งานบนเส้นทางวิกฤตไม่ขยับ \newline \textbf{แนวทาง} ในการประชุมประจำวัน ให้รายงานเป็นข้อแรกว่า ``วันนี้ทำอะไรที่ปลดล็อกคนอื่น'' ก่อนรายงานงานของตัวเอง & PM \newline เริ่มทันที \\
\bottomrule
\end{tabularx}

% =============================================================================
\section{การประเมินความเสี่ยง}

\begin{tabularx}{\dimexpr\textwidth-\leftskip\relax}{@{} L{3.8cm} L{1.4cm} L{1.6cm} Y @{}}
\toprule
\hrow \thead{ความเสี่ยง} & \thead{โอกาส} & \thead{ผลกระทบ} & \thead{แนวทางรับมือ} \\
\midrule
Sprint 2 แบกงานสอง Sprint รวมกันแล้วตามไม่ทัน & \textbf{สูง} & \textbf{สูง} & กำหนดจุดตัดสินใจกลาง Sprint ในวันที่ 20 ส.ค. หากวงจรหลักยังไม่ทะลุถึงขั้นสร้างคำขอได้ ให้ตัดหน้าจอที่ยังไม่ทำและคงเฉพาะเส้นทางที่สาธิตได้ \\
เชื่อมต่อครั้งแรกแล้วพบว่าโครงสร้างข้อมูลไม่ตรงกัน & ปานกลาง & สูง & ความเสี่ยงนี้ \textbf{ลดลง} จากครั้งที่ 1 เพราะงานทดลองสถาปัตยกรรมและเอกสารเตรียมสัญญาได้ระบุจุดที่ไม่ตรงกันไว้ล่วงหน้าแล้ว เหลือเพียงลงมือแก้ตามที่ระบุ \\
งานกระจุกที่ Sprint 3 ซึ่งตรงกับสอบกลางภาค & \textbf{สูง} & \textbf{สูง} & ความเสี่ยงนี้ \textbf{รุนแรงขึ้น} เพราะการเลื่อนหนึ่งชั้นทำให้งานกฎธุรกิจทั้งชุดมาตกที่ Sprint 3 พอดี ให้ย้ายงานที่ไม่ขึ้นกับผู้อื่นขึ้นมาทำใน Sprint 2 ให้มากที่สุด \\
ฐานข้อมูลขาดตารางผูกชิ้นและสถานะอนุมัติบางส่วน & \textbf{สูง} & ปานกลาง & ยังไม่ได้แก้มาสองรอบรายงาน ให้กำหนดเป็นงานบังคับของ Sprint 2 ขณะที่ฐานข้อมูลยังว่างและต้นทุนการแก้ยังต่ำ \\
กำลังคนฝั่ง Back-End ไม่พอเทียบสัดส่วนงาน & \textbf{สูง} & \textbf{สูง} & ความเสี่ยงนี้ \textbf{รุนแรงขึ้น} เพราะยืนยันด้วยหลักฐานแล้วว่าไม่มีความคืบหน้าเลย 5 วัน ให้ย้ายสมาชิกฝั่ง Front-End ที่ว่างจากการทำหน้าจอผู้ดูแลระบบเสร็จแล้วมาช่วยเขียน Procedure ที่ไม่ซับซ้อน \\
ทีมทำงานที่ทำได้ทันทีแทนงานที่ต้องรอ & \textbf{สูง} & ปานกลาง & ความเสี่ยงใหม่ที่พบในรอบนี้ ให้ PM ตรวจทุกวันว่ามีใครกำลังถูกบล็อกอยู่ และให้งานที่ปลดล็อกคนอื่นได้รับความสำคัญก่อนเสมอ \\
\bottomrule
\end{tabularx}

% =============================================================================
\clearpage
\section{แผนการดำเนินงานที่ปรับปรุงแล้ว}
\label{sec:replan}

หัวข้อนี้คือสาระสำคัญที่สุดของรายงานฉบับนี้ คณะทำงานเสนอ\textbf{เลื่อนแผนทั้งชุดไปหนึ่ง Sprint} โดยยึดวันที่ในปฏิทินเดิมไว้ไม่เปลี่ยน แต่เปลี่ยนเนื้อหาของแต่ละ Sprint ตามความเป็นจริง ตารางเทียบรายการทีละบรรทัดอยู่ในภาคผนวก ข

\subsection{ช่วงที่เหลือของ Sprint 1 (11 – 13 สิงหาคม 2569)}

\begin{tabularx}{\dimexpr\textwidth-\leftskip\relax}{@{} L{1.8cm} Y L{2.0cm} @{}}
\toprule
\hrow \thead{กำหนด} & \thead{งาน} & \thead{ผู้รับผิดชอบ} \\
\midrule
11 ส.ค. & เริ่มสคริปต์ข้อมูลตั้งต้น 2 หน่วยงาน พร้อมผู้ใช้ครบทุกบทบาท (ไม่ต้องรอสัญญา) & BE-dev \\
12 ส.ค. & เพิ่มตารางผูกชิ้นอุปกรณ์ สถานะอนุมัติบางส่วน และคีย์ไม่ซ้ำบนอีเมล เข้าสคีมา & BE-lead + SA \\
13 ส.ค. & \textbf{ประชุมลงมติสัญญา tRPC} แล้วคอมมิตไฟล์สัญญาเข้ารีโปในวันเดียวกัน & API + SA \\
13 ส.ค. & รวม CI เข้าสายหลัก & FE-lead \\
13 ส.ค. & เกณฑ์การยอมรับของวงจรหลัก & QA-lead \\
\bottomrule
\end{tabularx}

\vspace{0.8em}
\noindent\textbf{เกณฑ์ ``เสร็จ'' ที่ปรับใหม่ของ Sprint 1} สัญญาถูกล็อกและคอมมิตแล้ว · ฐานข้อมูลมีข้อมูลตั้งต้น · CI ทำงานบนสายหลัก \textbf{ส่วนงานวงจรยืม–คืนทั้งหมดเลื่อนไป Sprint 2}

\subsection{Sprint 2 ``ยืม–คืนครบวงจร'' (14 – 27 สิงหาคม 2569)}

Sprint นี้รับเป้าหมายเดิมของ Sprint 1 มาทั้งชุด บวกงานตกค้างจาก Sprint 0

\begin{itemize}
  \item \textbf{ฝั่ง Back-End} ยืนยันตัวตนและควบคุมสิทธิ์ตามสมาชิกต่อหน่วยงาน · รายการอุปกรณ์และจำนวนคงเหลือ · สร้างคำขอ อนุมัติอัตโนมัติ ผูกชิ้น ส่งมอบ และรับคืน · การจัดการข้อผิดพลาดและการตรวจความถูกต้อง
  \item \textbf{ฝั่ง Front-End} ยกระบบยืนยันตัวตนจริงเข้าแทนของจำลอง · ยกระดับไลบรารีตรวจสอบข้อมูลก่อนลงมือทำหน้าที่มีแบบฟอร์ม · หน้าเรียกดูอุปกรณ์ หน้าส่งคำขอ และหน้าคิวของเจ้าหน้าที่ · \textbf{เปลี่ยนหน้าจอผู้ดูแลระบบ 6 หน้าจากข้อมูลจำลองเป็นข้อมูลจริง}
  \item \textbf{ฝั่ง QA} แผนการทดสอบ เกณฑ์การยอมรับ และการทดสอบเส้นทางปกติรวมกรณียืมข้ามภาควิชา
\end{itemize}

\noindent\textbf{เกณฑ์ ``เสร็จ''} ยืมอุปกรณ์ของภาควิชาคอมพิวเตอร์และภาควิชาไฟฟ้าในระบบเดียวได้ครบวงจรในเส้นทางปกติ (ยังไม่มีรูปภาพและเครดิต) และ\textbf{ไม่เหลือข้อมูลจำลองในเส้นทางหลัก}

\noindent\textbf{ข้อควรระวัง} Sprint นี้แบกงานสอง Sprint รวมกัน คณะทำงานเสนอจุดตัดสินใจกลาง Sprint ในวันที่ 20 สิงหาคม หากถึงวันนั้นแล้วยังสร้างคำขอผ่านระบบจริงไม่ได้ ให้ตัดหน้าจอที่ยังไม่ได้ทำออกและคงไว้เฉพาะเส้นทางเดียวที่สาธิตได้

\subsection{Sprint 3 ``กฎธุรกิจ + หลายหน่วยงานระดับที่ 1'' (28 สิงหาคม – 10 กันยายน 2569)}

รับเนื้อหาเดิมของ Sprint 2 มาทั้งชุด ได้แก่ การอนุมัติตามระดับอุปกรณ์ การจัดเส้นทางอนุมัติตามหน่วยงานเจ้าของ การเลือกชิ้นอุปกรณ์ การจองชิ้นแบบอะตอมมิก การตรวจสภาพพร้อมรูปภาพ การประเมินความเสียหายและระบบเครดิต การแจ้งเตือน และเอกสาร API พร้อมหน้าจอฝั่งผู้ใช้ที่เกี่ยวข้องทั้งสี่หน้า

\noindent\textbf{ข้อควรระวัง} Sprint นี้ตรงกับช่วงสอบกลางภาคพอดี และหลังการเลื่อนหนึ่งชั้นแล้วกลายเป็น Sprint ที่มีงานหนาแน่นที่สุด ให้ย้ายงานที่ไม่ขึ้นกับผู้อื่นขึ้นไปทำใน Sprint 2 ให้มากที่สุดเท่าที่ทำได้

\subsection{Sprint 4 ``ตะกร้าข้ามหน่วยงาน + จองแบบย่อ + ขัดเงา'' (11 – 24 กันยายน 2569)}

รับเนื้อหาเดิมของ Sprint 3 มาในรูปแบบที่ตัดทอนแล้ว บวกงานขัดเงาเดิมของ Sprint 4

\begin{itemize}
  \item \textbf{คงไว้} ตะกร้าข้ามหน่วยงานพร้อมการแตกอนุมัติตามเจ้าของ · การจองล่วงหน้าแบบย่อ (จองตามชนิดอุปกรณ์และช่วงเวลา พร้อมการคำนวณจำนวนคงเหลือในอนาคต) · การต่ออายุพื้นฐาน · การบังคับสิทธิ์ทุกจุดพร้อมบันทึกการตรวจสอบ · การค้นหา การแบ่งหน้า การรองรับมือถือ และหน้าจอที่ยังเป็นหน้าเปล่า
  \item \textbf{ตัดออก} การจองสถานที่ · การส่งออกรายงาน · การจองล่วงหน้าขั้นสูง (ตามหัวข้อ \ref{sec:scopechange})
\end{itemize}

\noindent\textbf{ข้อควรระวัง} Sprint นี้เป็น\textbf{Sprint สุดท้ายที่เพิ่มความสามารถใหม่ได้} งานใดที่ไม่เสร็จภายในวันที่ 24 กันยายน ให้ถือว่าตัดออกจากขอบเขตทันที ห้ามลากเข้า Sprint 5

\subsection{Sprint 5 ``หยุดเพิ่มของและซ้อมนำเสนอ'' (25 กันยายน – 1 ตุลาคม 2569)}

ไม่มีการเปลี่ยนแปลงจากแผนเดิม ได้แก่ การหยุดเพิ่มความสามารถตั้งแต่ต้นสัปดาห์ การทดสอบถดถอยเต็มรูปแบบ การทดสอบสมรรถนะ การนำระบบขึ้นใช้งานจริง คู่มือผู้ใช้และคู่มือติดตั้ง และการซ้อมนำเสนอ

\subsection{การประเมินกำหนดเสร็จของโครงการ}

คณะทำงานประเมินว่า \textbf{โครงการยังเสร็จได้ภายในวันที่ 1 ตุลาคม 2569 ตามกำหนดเดิม} แต่เงื่อนไขเข้มขึ้นกว่ารายงานครั้งที่ 1 กล่าวคือต้องเป็นไปตามนี้ครบทั้งสามข้อ

\begin{enumerate}
  \item สัญญา API ถูกล็อกและคอมมิตภายในวันที่ 13 สิงหาคม 2569
  \item วงจรยืม–คืนเส้นทางปกติทำงานได้จริงภายในวันที่ 27 สิงหาคม 2569 ซึ่งเป็นวันสิ้นสุด Sprint 2
  \item คณะกรรมการอนุมัติการตัดขอบเขตทั้งสามรายการในรอบนี้
\end{enumerate}

\noindent หากข้อใดข้อหนึ่งไม่เป็นไปตามนี้ คณะทำงานจะเสนอ\textbf{ลดขอบเขตของ MVP ลงเหลือเฉพาะการยืมภายในหน่วยงานเดียว} โดยตัดความสามารถข้ามหน่วยงานออก ซึ่งจะกระทบจุดขายหลักของโครงการโดยตรง จึงขอเสนอเป็นทางเลือกสุดท้ายเท่านั้น

% =============================================================================
\section{ทรัพยากรและการสนับสนุนที่ต้องการเพิ่มเติม}

\begin{enumerate}
  \item \textbf{กำลังคนฝั่ง Back-End} เป็นข้อที่เร่งด่วนที่สุดในรอบนี้ ขณะนี้มีผู้เขียนโค้ดฝั่ง Back-End จริงเพียงคนเดียว ขณะที่ฝั่ง Back-End รับผิดชอบน้ำหนักงาน 30\% ซึ่งสูงที่สุดในโครงการ ข้อเสนอนี้ยื่นไว้ตั้งแต่รายงานครั้งที่ 1 และยังไม่ได้รับการแก้ไข
  \item \textbf{ฐานข้อมูล PostgreSQL สำหรับใช้ร่วมกันในทีม} ปัจจุบันแต่ละคนใช้ฐานข้อมูลบนเครื่องตัวเอง ทำให้ทดสอบร่วมกันไม่ได้ ค่าใช้จ่ายอยู่ในวงเงิน 1{,}500 บาทที่ประเมินไว้แล้ว
  \item \textbf{มติของที่ประชุมเรื่องสัญญา API} ซึ่งเป็นสิ่งเดียวที่ขวางไม่ให้ Sprint 2 เริ่มได้ งานเตรียมทั้งหมดพร้อมแล้ว
  \item \textbf{คำแนะนำจากอาจารย์ที่ปรึกษาเรื่องการคงไว้ซึ่ง MongoDB} เป็นข้อค้างจากรายงานครั้งที่ 1 ที่ยังไม่มีข้อสรุป
\end{enumerate}

% =============================================================================
\section{ความสำเร็จที่ควรบันทึกไว้}

\begin{itemize}
  \item \textbf{หน้าจอผู้ดูแลระบบเสร็จก่อนกำหนดสี่สัปดาห์} 6 หน้าประมาณ 2{,}000 บรรทัด เป็นงานที่แผนเดิมวางไว้ใน Sprint 3 การมีอยู่แล้วทำให้ Sprint 4 เบาลงจริง และเป็นหน้าจอชุดที่พร้อมสาธิตที่สุดในขณะนี้
  \item \textbf{งานเตรียมสัญญา API มีคุณภาพสูงผิดปกติสำหรับโครงการขนาดนี้} เอกสาร 64 หน้าพร้อมวาระ 19 วาระ และโค้ดตัวอย่าง 28 ไฟล์ ทำให้การประชุมลงมติใช้เวลาสั้นลงมาก และลดโอกาสที่จะตัดสินใจผิดแล้วต้องรื้อ
  \item \textbf{ระบบตรวจสอบอัตโนมัติถูกเขียนขึ้นแล้ว} และถูกตรวจสอบล่วงหน้าแล้วว่าจะผ่านกับสายหลักปัจจุบัน ซึ่งเป็นรายละเอียดที่มักถูกมองข้ามและทำให้ CI ที่ตั้งใหม่มักแดงตั้งแต่วันแรกจนไม่มีใครสนใจ
  \item \textbf{ชุดส่วนประกอบหน้าจอ 19 ตัว} รวมตารางข้อมูลและกราฟ ถูกสร้างขึ้นระหว่างทำหน้าจอผู้ดูแลระบบ และจะถูกใช้ซ้ำในหน้าที่เหลืออีก 19 หน้า
\end{itemize}

% =============================================================================
\clearpage
\section*{ภาคผนวก ก — ที่มาของตัวเลขความก้าวหน้า}
\addcontentsline{toc}{section}{ภาคผนวก ก — ที่มาของตัวเลขความก้าวหน้า}

ตัวเลขในหัวข้อ \ref{sec:dimension} ประเมินจากหลักฐานที่ตรวจสอบได้ใน Repository ณ วันที่ \reportdate\ โดยไม่นับโฟลเดอร์ \texttt{spike-option-b/}

\begin{tabularx}{\dimexpr\textwidth-\leftskip\relax}{@{} L{2.9cm} L{1.1cm} Y @{}}
\toprule
\hrow \thead{มิติงาน} & \thead{ผลจริง} & \thead{หลักฐานและวิธีคำนวณ} \\
\midrule
Back-End Design & 55\% & ให้น้ำหนักแบบจำลองข้อมูลและสัญญา API อย่างละครึ่ง \newline \textbf{แบบจำลองข้อมูล 85\%} จาก \texttt{schema.prisma} ที่มีครบ 33 ตารางและ Migration \texttt{20260805125233\_init} ที่สร้างตารางได้จริง หักส่วนที่ยังขาดคือตารางผูกชิ้นอุปกรณ์ สถานะอนุมัติบางส่วน และคีย์ไม่ซ้ำ \newline \textbf{สัญญา API 25\%} เพิ่มจาก 5\% ในรอบที่แล้ว เพราะมีคู่มือ 31 หน้า วาระประชุม 33 หน้า และโค้ดตัวอย่าง 28 ไฟล์ แต่ยังไม่มีไฟล์สัญญาในรีโปและยังไม่มีมติ \newline $(85+25)/2 = 55$ \\
Front-End Design & 72\% & ระบบออกแบบและเปลือกระบบ (น้ำหนัก 40\%) เสร็จ 100\% · แบบหน้าจอ (น้ำหนัก 35\%) เสร็จ 7 จาก 26 หน้า คิดเป็น 27\% · การออกแบบกระบวนการทางธุรกิจ (น้ำหนัก 25\%) เสร็จ 90\% จากข้อเสนอโครงการ \newline $0.40(100)+0.35(27)+0.25(90) = 72$ \\
Back-End Dev & 12\% & \textbf{ไม่เปลี่ยนจากรอบที่แล้ว} คอมมิตล่าสุดของ \texttt{backend/} คือ \texttt{ed18534} วันที่ 5 ส.ค. นับเฉพาะโครงสร้างพื้นฐานที่ประกอบเสร็จใน \texttt{app.module.ts} และ Migration ที่รันได้ ไฟล์ใน \texttt{src/} มีเพียงโครงเริ่มต้นของ NestJS กับไคลเอนต์ Prisma ที่ถูกสร้างอัตโนมัติ ไม่มี Router ไม่มีการยืนยันตัวตน ไม่มี Procedure และไม่มีสคริปต์ข้อมูลตั้งต้น \\
Front-End Dev & 38\% & ให้น้ำหนักหน้าจอ 70\% และโครงสร้างพื้นฐาน 30\% \newline \textbf{หน้าจอ 19\%} จาก 26 หน้า มีเนื้อหาจริง 7 หน้า แต่ทุกหน้าอ่านจากข้อมูลจำลอง จึงคิดให้หน้าละ 0.7 ของหน้าที่ต่อข้อมูลจริงแล้ว ได้ $7\times0.7 = 4.9$ เทียบ 26 หน้า คิดเป็น 19\% \newline \textbf{โครงสร้างพื้นฐาน 82\%} เปลือกระบบ เมนูตามบทบาท ส่วนประกอบ 19 ตัว i18n ธีม และโครงยืนยันตัวตน เสร็จหมด หักส่วนที่ยังไม่มีผล คือ CI ที่ยังไม่รวมเข้าสายหลัก และ \texttt{api-client.ts} ที่ยังไม่มีผู้เรียกใช้ \newline $0.70(19)+0.30(82) = 38$ \\
Test & 0\% & \textbf{ไม่เปลี่ยนจากรอบที่แล้ว} ไม่พบไฟล์ทดสอบที่เขียนขึ้นเองทั้งใน \texttt{backend/src} และ \texttt{frontend/src} พบเพียง \texttt{app.controller.spec.ts} กับ \texttt{prisma.service.spec.ts} ซึ่งเป็นไฟล์ตัวอย่างที่มากับโครงเริ่มต้นของ NestJS และไม่พบเอกสารแผนการทดสอบ \\
Documentation & 62\% & เพิ่มจาก 45\% จากเอกสารใหม่ในรอบนี้ ได้แก่ คู่มือสัญญา 31 หน้า วาระประชุม 33 หน้า โค้ดตัวอย่าง 28 ไฟล์ และสารบัญไฟล์ของโครงการ รวมกับข้อเสนอโครงการ 42 หน้าและรายงานความก้าวหน้าครั้งที่ 1 ที่มีอยู่เดิม \newline ส่วนที่ยังขาดคือ SRS ฉบับทางการ SDS นอกเหนือบทฐานข้อมูล คู่มือผู้ใช้ และคู่มือติดตั้ง \\
\bottomrule
\end{tabularx}

\vspace{1em}
\noindent\textbf{หมายเหตุเรื่องตัวเลขสองชุด} เอกสารตรวจรายการ MVP ต่อ Sprint (ไฟล์ ``ULMs — MVP Checklist ต่อ Sprint (ปรับปรุง 10 ส.ค. 2569)'' นามสกุล \texttt{.xlsx} ในโฟลเดอร์ \texttt{docs/resource/}) แสดงความก้าวหน้ารวมที่ 24\% ซึ่งต่างจาก \actualpct\ ในรายงานฉบับนี้ ตัวเลขทั้งสองไม่ขัดแย้งกันเพราะวัดคนละอย่าง เอกสารตรวจรายการนับ\textbf{จำนวนงาน}โดยให้ทุกงานมีน้ำหนักเท่ากันและเฉลี่ยข้าม Sprint ทั้งหก ซึ่ง Sprint ที่ยังไม่เริ่มสี่ Sprint จะถ่วงค่าลงเป็นศูนย์ ส่วนรายงานฉบับนี้ถ่วงน้ำหนักตาม\textbf{ปริมาณงาน}ของแต่ละมิติตามที่ใช้มาตั้งแต่ครั้งที่ 1 จึงเทียบกับรายงานครั้งก่อนได้ตรง ๆ

\vspace{0.5em}
\noindent\textbf{ข้อจำกัดของการประเมิน} ตัวเลขข้างต้นเป็นการประเมินจากปริมาณผลงานที่ปรากฏใน Repository ซึ่งตรวจสอบย้อนกลับได้ทุกบรรทัด แต่การแปลงปริมาณผลงานเป็นร้อยละย่อมมีดุลพินิจของผู้ประเมินอยู่ด้วย โดยเฉพาะค่าน้ำหนักของแต่ละมิติและค่าปรับลด 0.7 ที่ใช้กับหน้าจอที่ยังอ่านจากข้อมูลจำลอง คณะทำงานจึงรายงานทั้งค่าน้ำหนักและวิธีคำนวณไว้อย่างเปิดเผย เพื่อให้ผู้อ่านปรับค่าตามที่เห็นสมควรได้

% =============================================================================
\clearpage
\section*{ภาคผนวก ข — ตารางเทียบแผนเดิมกับแผนที่ปรับปรุงแล้ว}
\addcontentsline{toc}{section}{ภาคผนวก ข — ตารางเทียบแผนเดิมกับแผนที่ปรับปรุงแล้ว}

ตารางนี้แจกแจงทุกรายการที่เปลี่ยนตำแหน่งหรือถูกตัด เพื่อให้ตรวจสอบได้ว่าไม่มีงานใดหายไปโดยไม่มีใครสังเกต

\begin{tabularx}{\dimexpr\textwidth-\leftskip\relax}{@{} L{1.6cm} Y L{1.8cm} L{1.8cm} @{}}
\toprule
\hrow \thead{ประเภท} & \thead{รายการ} & \thead{เดิม} & \thead{ใหม่} \\
\midrule
เลื่อน & ยืนยันตัวตนและควบคุมสิทธิ์ตามสมาชิกหน่วยงาน & Sprint 1 & Sprint 2 \\
เลื่อน & รายการอุปกรณ์ ค้นหา และจำนวนคงเหลือ & Sprint 1 & Sprint 2 \\
เลื่อน & สร้างคำขอ อนุมัติอัตโนมัติ ผูกชิ้น ส่งมอบ รับคืน & Sprint 1 & Sprint 2 \\
เลื่อน & หน้าเรียกดูอุปกรณ์ · หน้าส่งคำขอ · หน้าคิวเจ้าหน้าที่ & Sprint 1 & Sprint 2 \\
เลื่อน & ทดสอบเส้นทางปกติและกรณียืมข้ามภาควิชา & Sprint 1 & Sprint 2 \\
เลื่อน & ข้อมูลตั้งต้น 2 หน่วยงาน & Sprint 0 & Sprint 2 \\
เลื่อน & แผนการทดสอบและเกณฑ์การยอมรับ & Sprint 0 & Sprint 2 \\
\midrule
เลื่อน & อนุมัติตามระดับอุปกรณ์ (T0/T1/T2) & Sprint 2 & Sprint 3 \\
เลื่อน & จัดเส้นทางอนุมัติตามหน่วยงานเจ้าของ & Sprint 2 & Sprint 3 \\
เลื่อน & เลือกชิ้นอุปกรณ์ · จองชิ้นแบบอะตอมมิก & Sprint 2 & Sprint 3 \\
เลื่อน & ตรวจสภาพพร้อมรูปภาพ · ประเมินความเสียหาย · เครดิต & Sprint 2 & Sprint 3 \\
เลื่อน & การแจ้งเตือน · เอกสาร API & Sprint 2 & Sprint 3 \\
เลื่อน & หน้าจอรับของ · ประเมินความเสียหาย · อนุมัติ · เครดิต & Sprint 2 & Sprint 3 \\
\midrule
เลื่อน & ตะกร้าข้ามหน่วยงานและสถานะอนุมัติบางส่วน & Sprint 3 & Sprint 4 \\
เลื่อน & การจองล่วงหน้า (เฉพาะส่วนแกน) & Sprint 3 & Sprint 4 \\
เลื่อน & บังคับสิทธิ์ทุกจุด · แยกหน้าที่ · บันทึกการตรวจสอบ & Sprint 3 & Sprint 4 \\
\midrule
ตัด & การจองสถานที่ (Facility T3) ทั้งฝั่ง Back-End และหน้าจอ & Sprint 4 & \textbf{ตัดออก} \\
ตัด & การส่งออกรายงาน & Sprint 4 & \textbf{ตัดออก} \\
ตัด & ผูกชิ้นอัตโนมัติ · จัดคิวเมื่อจองไม่สำเร็จ · เพดานต่ออายุเทียบวันจอง & Sprint 3 & \textbf{ตัดออก} \\
\midrule
ดึงขึ้น & หน้าจอผู้ดูแลระบบ 6 หน้า & Sprint 3 & \textbf{เสร็จแล้ว} \\
\midrule
เพิ่ม & ระบบตรวจสอบอัตโนมัติ CI 3 งานตรวจ & ไม่มีในแผน & Sprint 1 \\
เพิ่ม & เอกสารเตรียมสัญญา tRPC 64 หน้า และโค้ดตัวอย่าง 28 ไฟล์ & ไม่มีในแผน & Sprint 1 \\
เพิ่ม & เพิ่มตารางผูกชิ้น สถานะอนุมัติบางส่วน และคีย์ไม่ซ้ำ เข้าสคีมา & ไม่มีในแผน & Sprint 2 \\
เพิ่ม & ยกระบบยืนยันตัวตนจริงเข้าแทนของจำลอง และยกระดับไลบรารีตรวจข้อมูล & ไม่มีในแผน & Sprint 2 \\
เพิ่ม & เปลี่ยนหน้าจอผู้ดูแลระบบจากข้อมูลจำลองเป็นข้อมูลจริง & ไม่มีในแผน & Sprint 2 \\
\bottomrule
\end{tabularx}

\vspace{1em}
\noindent สรุปเป็นตัวเลข \textbf{เลื่อน 16 รายการ · ตัด 3 รายการ · ดึงขึ้นมาทำก่อน 1 รายการ · เพิ่มใหม่ 5 รายการ} โดยรายการที่เพิ่มใหม่ทั้งห้าเป็นงานที่จำเป็นอยู่แล้วแต่ไม่เคยปรากฏในแผนฉบับเดิม จึงไม่เคยถูกนับเป็นภาระงานของ Sprint ใด

\end{document}
